% !TeX root = ../../Book2.tex
% 纲要标题：## 第2章　模、子模与商模

\chapter{模、子模与商模}
\label{b2:ch:02}
\BookChapterNumber{b2:ch:02}

\begin{chapterabstract}
把向量空间的标量域换成环，便得到模。整数模统一了交换群的运算，多项式环上的模则把线性算子包含在标量作用中。本章建立子模、商模、核、像与余核，证明模的同构定理，再用生成元、零化子和挠性辨认模的特点。最后构造任意族的直和与直积，说明它们在有限情形下一致，而在无限情形下承担不同的映射任务。
\end{chapterabstract}

\begin{learninggoals}
\begin{itemize}
\item 能为交换群、矩阵作用和线性算子写出相应的模结构，并判断同态与子模。
\item 会构造商模，求具体同态的核、像与余核，写出同构定理中的实际映射。
\item 区分单个元素与整个模的零化子，理解忠实、有限生成、挠与无挠各自的含义。
\item 掌握直和、直积的有限支集条件及泛性质，判断无限情形下的求和与坐标论证是否有效。
\end{itemize}
\end{learninggoals}

整数可以作用在任何交换群上：\(n\) 的作用就是重复相加或取加法逆元。以 \(\mathbb Z/6\mathbb Z\) 为例，一个元素 \(\bar 1\) 生成了整个群，却满足 \(6\bar 1=0\)。它有生成元，却不能像向量空间那样用一组基作无关系的线性组合。若把标量从整数换成一般环，生成元、关系、商对象和同态仍然有意义；哪些向量空间中的推理能够保留，则要由环上的乘法来判断。

\ProseParagraphBreak

给定一个子模，商群上的标量作用怎样定义，才不会依赖代表元？同态的核与像又怎样决定相应的商？遇到无限族模时，一个元组可以有多少个非零坐标？这些问题分别引出商模、同构定理以及直和与直积。构造它们不必要求非零标量可逆；但选基和拆出补模等向量空间里的便利，在一般环上不能照搬。

\ProseParagraphBreak

本章需要交换群、环、左理想与双边理想的定义，以及第一章的向量空间和线性映射。商模与同构定理同时适用于交换环和非交换环。非交换环上须区分左作用与右作用，同态组成的集合也不再自动具有原环上的自然模结构；定义模同态后，我们会说明增加标量作用所需的条件。整环上的挠性另在\subsecref{b2:subsec:0203:03}集中讨论，不能未经说明推广到任意环。

\ProseParagraphBreak
模、同态和直和的基础理论可参阅 Lang 的《Algebra》\cite[CHAPTER III, \S\S 1--3]{Lang2002}。Roman 的《Advanced Linear Algebra》\cite[Chapter 4, pp. 109--124]{Roman2008AdvancedLinearAlgebra}从线性代数转入模论，适合对照本章的基本构造；该部分以交换环为范围，本章对非交换环的左右作用另作说明。

% !TeX root = ../../Book2.tex
% 纲要标题：### 2.1 环上的模

\section{环上的模}
\label{b2:sec:0201}

设 \(k\) 是域，\(A,B\in M_n(k)\)，\(v\) 与 \(w\) 分别是长度为 \(n\) 的列向量和行向量。先用 \(A\)、再用 \(B\) 作用于 \(v\)，得到 \(B(Av)=(BA)v\)；同样作用于 \(w\)，则得到 \((wA)B=w(AB)\)。前者对应左模的结合律 \((rs)m=r(sm)\)：先作用 \(s\)，再作用 \(r\)，合起来由 \(rs\) 作用；后者对应右模的结合律 \((mr)s=m(rs)\)：先作用 \(r\)，再作用 \(s\)，合起来也由 \(rs\) 作用。标量写在哪一侧，记录的是环乘法与作用复合的次序。

% 纲要要点（供目录核对）：
% * 左模、右模与标量约定
% * 交换群、向量空间及线性算子模
% * 正则模与表示

\subsection{左模、右模与标量约定}
\label{b2:subsec:0201:01}

本章的环均含单位元，允许非交换环和零环，环同态保持单位元。除特别说明外，模指幺性左模。环的零元与模的零元通常都写成 \(0\)，必要时分别写 \(0_R\) 与 \(0_M\)。

\begin{definition}\label{b2:def:left-module}
设 \(R\) 是环，\((M,+)\) 是交换群，并给定映射
\[
R\times M\longrightarrow M,\qquad (r,m)\longmapsto rm,
\]
这里 \(rm\) 也可写为 \(r\cdot m\)，两者表示这一给定的左标量乘法。
如果任意 \(r,s\in R\) 和 \(m,n\in M\) 都满足
\[
r(m+n)=rm+rn,\qquad (r+s)m=rm+sm,
\]
\[
(rs)m=r(sm),\qquad 1_Rm=m.
\]
就称 \(M\) 为\term[zuomo]{左 \(R\) 模}[left R-module]。最后一条称为幺性条件。这里的加法交换群已经包括加法结合律、零元、负元与交换律。
\end{definition}

由分配律，\(0_Rm=(0_R+0_R)m=0_Rm+0_Rm\)，在交换群中消去一项便得 \(0_Rm=0_M\)。同样，\(r0_M=0_M\)，且
\[
(-r)m=-(rm)=r(-m).
\]
这些恒等式可以从公理推出，无须另列为模的定义。若 \(R\) 是零环，则 \(1_R=0_R\)，所以 \(m=1_Rm=0_Rm=0_M\)；零环只有零模。对任何环，只有一个元素的交换群都有唯一的模结构，称为\term[lingmo]{零模}[zero module]，也记为 \(0\)。

\begin{definition}\label{b2:def:right-module}
设 \(R\) 是环，\((M,+)\) 是交换群，并给定右侧标量乘法 \(M\times R\to M\)、\((m,r)\mapsto mr\)。这里 \(mr\) 也可写为 \(m\cdot r\)，两者表示这一给定的右标量乘法。如果对任意 \(m,n\in M\) 和 \(r,s\in R\) 都有
\[
(m+n)r=mr+nr,\qquad m(r+s)=mr+ms,
\]
\[
(mr)s=m(rs),\qquad m1_R=m.
\]
就称 \(M\) 为\term[youmo]{右 \(R\) 模}[right R-module]。
\end{definition}

若 \(R\) 交换，一个左模可以通过 \(mr=rm\) 看成右模，因为 \((mr)s=s(rm)=(sr)m=(rs)m\)。非交换时，这一步中的 \(sr=rs\) 一般不成立。例如，\(M_n(k)\) 通过左乘作用在列向量 \(k^n\) 上，而通过右乘作用在行向量上；在同一个列向量空间中把 \(vA\) 定义为 \(Av\)，所得右作用需要 \(BA v=AB v\)，通常不满足。

\ProseParagraphBreak
为了把左模的定义和定理用于右模，我们可以保留被作用的交换群，只把标量环的乘法次序倒过来。这样建立左模理论后便能得到相应的右模理论，无须重复整套证明。设 \(R\) 是环，在同一个加法群上定义新乘法 \(r*s=sr\)。所得环称为 \(R\) 的\term[fanhuan]{反环}[opposite ring]，记为 \(R^{\mathrm{op}}\)。一个右 \(R\) 模成为左 \(R^{\mathrm{op}}\) 模的方法是 \(r\cdot m=mr\)：确有
\[
(r*s)\cdot m=m(sr)=(ms)r=r\cdot(s\cdot m).
\]
反过来，从左 \(R^{\mathrm{op}}\) 模按 \(mr=r\cdot m\) 定义作用也得到右 \(R\) 模，且这两个转换互逆。
\symidx{\(R^{\mathrm{op}}\)：环 \(R\) 的反环}

\begin{table}[htbp]
\centering
\caption{左模、右模与反环的作用次序}\label{b2:tab:module-sides}
\begin{tabularx}{\linewidth}{@{}p{0.26\linewidth}p{0.22\linewidth}X@{}}
\toprule
结构 & 标量写法 & 结合律 \\
\midrule
左 \(R\) 模 & \(rm\) & \((rs)m=r(sm)\) \\
右 \(R\) 模 & \(mr\) & \((mr)s=m(rs)\) \\
左 \(R^{\mathrm{op}}\) 模，由右 \(R\) 模得到 & \(r\cdot m=mr\) & \((r*s)\cdot m=r\cdot(s\cdot m)\)，其中 \(r*s=sr\) \\
\bottomrule
\end{tabularx}
\end{table}
在\cref{b2:tab:module-sides}的最后一行，两边按右 \(R\) 作用计算均为 \(m(sr)\)；反环的乘法恰好补偿了作用方向的改变。

\subsection{交换群、向量空间及线性算子模}
\label{b2:subsec:0201:02}

任意交换群 \(A\) 都有唯一的 \(\mathbb Z\) 模结构。对正整数 \(n\)，令 \(na\) 为 \(n\) 个 \(a\) 的和；令 \(0a=0\) 与 \((-n)a=-(na)\)。整数乘法的分配律和结合律给出模公理。唯一性来自 \(1a=a\) 及对标量加法的分配律：它们已经决定了每个整数的作用。因而，交换群的整数倍运算就是最初的模运算。

\ProseParagraphBreak
域 \(k\) 上的模恰好是 \(k\) 向量空间。相比之下，\(\mathbb Z/6\mathbb Z\) 作为 \(\mathbb Z\) 模具有非零元素 \(\bar2\)，却有 \(3\bar2=0\)。在向量空间中，\(a\ne0\) 时可以乘 \(a^{-1}\)，由 \(av=0\) 推出 \(v=0\)；这里的整数 \(3\) 没有逆元，这个推理便不能使用。模的研究首先要求把这种差别与熟悉的线性运算区别开来。

\ProseParagraphBreak
\secref{b2:sec:0103}已经用多项式作用研究线性算子。下面把同一构造写成正式的模结构，并说明如何从模的作用恢复算子。

\begin{proposition}\label{b2:prop:operator-module}
设 \(k\) 为域，\(V\) 为一个固定的 \(k\) 向量空间。\(V\) 上保持原有加法、延拓原有 \(k\) 标量作用的幺性左 \(k[t]\) 模结构，与 \(\operatorname{End}_k(V)\) 中的线性算子一一对应。给定 \(T\in\operatorname{End}_k(V)\)，对 \(f(t)=\sum_{j=0}^d a_jt^j\in k[t]\)，相应作用为
\[
f(t)v=f(T)v=\sum_{j=0}^d a_jT^j(v),\qquad T^0=\operatorname{id}_V.
\]
所得模记为 \(V_T\)。反向对应把模的作用送到线性算子 \(T(v)=tv\)。
\end{proposition}
\begin{proof}
每个 \(T^j\) 都线性，所以关于向量和与多项式和的两条分配律成立。若 \(g(t)=\sum_i b_it^i\)，则有限求和可以按指数重排，得到
\[
f(T)g(T)v=\sum_{j,i}a_jb_iT^{j+i}(v)=(fg)(T)v.
\]
常数多项式 \(1\) 作用为恒等映射，故模公理成立；常数多项式 \(a\) 的作用为原有的标量乘法 \(v\mapsto av\)。

反过来，给定延拓原有 \(k\) 作用的左 \(k[t]\) 模结构，定义 \(T(v)=tv\)。分配律给出可加性；对 \(a\in k\)，由 \(ta=at\) 得 \(T(av)=a\cdot T(v)\)，所以 \(T\) 对原有的 \(k\) 向量空间结构线性。再用乘法相容性归纳得到 \(t^jv=T^j(v)\)，由有限和的分配律得到所有多项式的作用公式。从算子构造模再取 \(t\) 的作用，恢复的正是原算子；从模的作用取出算子再计算 \(f(T)\)，恢复的正是原多项式作用，故两个构造互逆。
\end{proof}
\symidx{\(V_T\)：线性算子 \(T\) 确定的 \(k[t]\) 模}

任意左 \(k[t]\) 模限制到常数子环 \(k\) 后都成为 \(k\) 向量空间；在这个向量空间上，上述对应适用。

\ProseParagraphBreak
例如，取 \(V=k^2\)，令 \(T(a,b)=(b,0)\)。由于 \(T^2=0\)，任意多项式 \(f(t)=a_0+a_1t+t^2\cdot g(t)\) 的作用简化为
\[
f(t)(a,b)=(a_0a+a_1b,a_0b).
\]
这个模中的 \(t\) 既不是零作用，也不是可逆作用。同一个向量空间若改用零算子，\(t\) 就零化所有向量。仅知道底层向量空间的维数，尚不能辨认这两种多项式环上的模；它们的差异正是算子的差异。下一节定义模同态后，\cref{b2:prop:operator-module-intertwiners}将说明，保持多项式作用的同构恰好是交织两个算子的线性同构，因而这正对应第一章的相似关系。

\subsection{正则模与表示}
\label{b2:subsec:0201:03}

环可以作用在自身上。以环的加法群为底群，以环乘法为左标量作用，得到\term[zuozhengze-mo]{左正则模}[left regular module] \({}_RR\)；右乘则给出\term[youzhengze-mo]{右正则模}[right regular module] \(R_R\)。两者的底集合相同，标量作用的方向不同。更一般地，环同态 \(\varphi:R\to S\) 可以把任意左 \(S\) 模 \(M\) 变成左 \(R\) 模，规定
\[
r\cdot m=\varphi(r)m.
\]
例如 \(\varphi(rs)=\varphi(r)\varphi(s)\) 保证乘法相容，\(\varphi(1_R)=1_S\) 保证幺性。这称为沿 \(\varphi\) 的\term[xianzhi-biaoliang]{限制标量}[restriction of scalars]。
沿包含映射 \(\mathbb R\hookrightarrow\mathbb C\) 限制标量，一维复向量空间 \(\mathbb C\) 成为以 \(1,i\) 为基的二维实向量空间：元素与加法不变，只允许实数作标量。沿商映射 \(\mathbb Z\to\mathbb Z/n\mathbb Z\)（\(n\ge1\)）限制标量，则把一个 \(\mathbb Z/n\mathbb Z\) 模看成交换群上的整数倍作用，整数 \(r\) 的作用由其剩余类 \(\bar r\) 给出。

\ProseParagraphBreak
从映射的角度看，一个模给每个标量指定了底群的一个自同态。为准确说明这句话，设 \(A\) 为交换群，记 \(\operatorname{End}_{\mathbb Z}(A)\) 为其所有加法群自同态组成的集合。它具有逐点加法与复合乘法，约定 \(uv=u\circ v\)，即先作用 \(v\)，再作用 \(u\)。两条分配律分别来自逐点加法的定义与 \(u\) 的可加性；加法负元逐点取负，乘法单位为恒等映射。因此它是一个环。
\symidx{\(\operatorname{End}_{\mathbb Z}(A)\)：交换群 \(A\) 的自同态环}

\begin{proposition}\label{b2:prop:module-representation}
设 \(R\) 是环，\(A\) 是交换群。给 \(A\) 一个幺性左 \(R\) 模结构，等价于给一个保幺环同态
\[
\rho:R\longrightarrow\operatorname{End}_{\mathbb Z}(A).
\]
对应关系为 \(\rho(r)(a)=ra\)。这个环同态称为环 \(R\) 在 \(A\) 上的一个\term[huan-biaoshi]{表示}[representation of a ring]。
\end{proposition}
\begin{proof}
给定模结构，关于加法群运算的分配律说明每个 \(\rho(r)\) 都是加法自同态。其余模公理依次给出
\[
\rho(r+s)=\rho(r)+\rho(s),\qquad
\rho(rs)=\rho(r)\circ\rho(s),\qquad
\rho(1_R)=\operatorname{id}_A.
\]
所以 \(\rho\) 是保幺环同态。反过来，用 \(ra=\rho(r)(a)\) 定义作用，自同态的可加性给出第一条分配律，同态的三项性质给出另外三条模公理。两种构造逐个标量、逐个元素地互相恢复。
\end{proof}

对 \(M_n(k)\) 在列向量 \(k^n\) 上的作用，\(\rho(A)\) 就是左乘矩阵 \(A\) 的线性变换。对 \(V_T\)，表示把多项式 \(f\) 送到 \(f(T)\)。对左正则模，\(\rho(r)\) 是左乘 \(r\) 的映射，而且从 \(\rho(r)(1_R)=r\) 可以恢复标量本身。模与表示在这里描述的是同一份数据：前者突出被作用的对象，后者突出标量所确定的映射。

% !TeX root = ../../Book2.tex
% 纲要标题：### 2.2 子模、商模与模同态

\section{子模、商模与模同态}
\label{b2:sec:0202}

模之间的映射若同时保持加法和标量作用，它的核与像便能像在线性代数中那样用来分析映射。要进一步把核压成零，或把像从陪域中模掉，还须知道子模怎样形成商模。这些构造同样适用于非交换环上的左模。

% 纲要要点（供目录核对）：
% * 子模、商模及其泛性质
% * 核、像、余核与同态群；算子交织与自同态环
% * 模的同构定理

\subsection{子模、商模及其泛性质}
\label{b2:subsec:0202:01}

\begin{definition}\label{b2:def:module-homomorphism}
设 \(R\) 是环，\(M,N\) 是左 \(R\) 模。若映射 \(f:M\to N\) 对所有 \(m,m'\in M\) 及 \(r\in R\) 都满足
\[
f(m+m')=f(m)+f(m'),\qquad f(rm)=r\cdot f(m).
\]
就称 \(f\) 为 \(R\) \term[mo-tongtai]{模同态}[module homomorphism]或 \(R\) 线性映射。若模同态是双射，就称它为\term[mo-tonggou]{模同构}[module isomorphism]；存在这种同构时记 \(M\cong N\)。
\end{definition}

同态自动把零元送到零元、把负元送到负元。若 \(f\) 双射，对 \(n=f(m)\)、\(n'=f(m')\)，有 \(f^{-1}(n+n')=m+m'\) 与 \(f^{-1}(rn)=rm\)，故逆映射也是模同态。\(\mathbb Z\) 模同态正是交换群同态，因为可加性迫使它保持全部整数倍。

\begin{proposition}\label{b2:prop:operator-module-intertwiners}
设 \(V,W\) 是域 \(k\) 上的向量空间，\(T\in\operatorname{End}_k(V)\)、\(S\in\operatorname{End}_k(W)\)，相应的算子模为 \(V_T,W_S\)。映射 \(F:V\to W\) 是 \(k[t]\) 模同态，当且仅当它 \(k\) 线性且满足 \(FT=SF\)。因此 \(V_T\cong W_S\) 当且仅当存在满足此等式的 \(k\) 线性同构；在同一个有限维空间上，这等价于两个算子的矩阵相似。
\end{proposition}
\begin{proof}
模同态保持常数多项式的作用，故 \(k\) 线性；保持 \(t\) 的作用恰好给出 \(FT=SF\)。反过来，\cref{b2:prop:polynomial-action-subspace}说明 \(k\) 线性映射满足 \(FT=SF\) 时，对每个 \(p\in k[t]\) 都有 \(Fp(T)=p(S)F\)，所以它保持全部标量作用，是模同态。可逆时等式改写为 \(S=FTF^{-1}\)，得到相似关系。
\end{proof}

由此，第一章的算子相似问题就成为算子模的同构问题；所比较的既是向量空间，也是指定的多项式作用。

\begin{definition}\label{b2:def:submodule}
设 \(R\) 是环，\(M\) 是左 \(R\) 模。若子集 \(L\subseteq M\) 包含 \(0\)，且对任意 \(x,y\in L\)、\(r\in R\) 都有 \(x+y\in L\)、\(-x\in L\) 与 \(rx\in L\)，就称 \(L\) 为 \(M\) 的\term[zimo]{子模}[submodule]。它在限制的运算下本身是左 \(R\) 模，记为 \(L\le M\)。
\end{definition}
\symidx{\(L\le M\)：子模包含记号}

一个非空子集是子模，当且仅当它对 \(x-y\) 与 \(rx\) 封闭：取 \(x=y\) 得零元，取 \(x=0\) 得负元，再由差得到和。左正则模 \({}_RR\) 的子模恰为左理想；不能把这里的左理想自动当作双边理想。对于算子模 \(V_T\)，子模恰为满足 \(T(W)\subseteq W\) 的 \(k\) 线性子空间 \(W\)：一方面它必须对常数与 \(t\) 的作用封闭，另一方面这种封闭性使它对每个 \(T^j\) 和每个 \(f(T)\) 都封闭。

\begin{proposition}\label{b2:prop:quotient-module}
设 \(R\) 是环，\(M\) 是左 \(R\) 模，\(L\le M\)。对 \(m,m'\in M\)、\(r\in R\)，在陪集集合 \(M/L=\{\,m+L\mid m\in M\,\}\) 上定义
\[
(m+L)+(m'+L)=(m+m')+L,\qquad r(m+L)=rm+L.
\]
这给出左 \(R\) 模，称为\term[shangmo]{商模}[quotient module]。商映射 \(q:M\to M/L\)、\(q(m)=m+L\) 是满同态。对任意左 \(R\) 模 \(P\) 及同态 \(g:M\to P\)，存在同态 \(\bar g:M/L\to P\) 满足 \(g=\bar gq\)，当且仅当 \(g(L)=0\)；存在时 \(\bar g\) 唯一。
当 \(g(L)=0\) 时，所要求的分解由交换图表示为
\[
\begin{tikzcd}[column sep=3.5em,row sep=2.2em]
M\arrow[r,"g"]\arrow[d,"q"']&P\\
M/L\arrow[ur,dashed,"\bar g"']&
\end{tikzcd}
\]
\end{proposition}
\symidx{\(M/L\)：模按子模所得的商模}
\begin{proof}
陪集相等 \(m+L=m'+L\) 等价于 \(m-m'\in L\)。如果另取代表元 \(m+\ell\)、\(m'+\ell'\)，两者之和与 \(m+m'\) 的差为 \(\ell+\ell'\in L\)，而 \(r(m+\ell)-rm=r\ell\in L\)。所以两项运算均不依赖代表元。加法的结合律与交换律由 \(M\) 中对应等式下降，零元是 \(L\)，\(m+L\) 的负元是 \(-m+L\)。例如标量乘法的相容性为
\[
(rs)(m+L)=(rs)m+L=r(sm)+L=r\bigl(s(m+L)\bigr),
\]
另外两条分配律和幺性同样由 \(M\) 中的公理得到。公式本身说明 \(q\) 是满同态。

若 \(g=\bar gq\)，则对 \(\ell\in L\)，有 \(g(\ell)=\bar g(0)=0\)。反过来，若 \(g(L)=0\)，定义 \(\bar g(m+L)=g(m)\)。当 \(m-m'\in L\) 时，\(g(m)-g(m')=g(m-m')=0\)，故定义良好。代入商模运算公式，便得到 \(\bar g\) 的可加性与 \(R\) 线性。由于每个陪集都有形式 \(q(m)\)，关系 \(g=\bar gq\) 迫使这个公式成立，也就给出唯一性。
\end{proof}

\subsection{核、像、余核与同态群}
\label{b2:subsec:0202:02}

\begin{definition}\label{b2:def:module-kernel-image-cokernel}
设 \(R\) 是环，\(f:M\to N\) 是左 \(R\) 模之间的同态。定义其\term[he]{核}[kernel]、\term[xiang]{像}[image]及\term[yuhe]{余核}[cokernel]为
\[
\ker f=\{\,m\in M\mid f(m)=0\,\},\qquad
\operatorname{im}f=\{\,f(m)\mid m\in M\,\},
\]
\[
\operatorname{coker}f=N/\operatorname{im}f.
\]
其中核是 \(M\) 的子模，像是 \(N\) 的子模，余核附有商映射 \(q_f:N\to\operatorname{coker}f\)。
\end{definition}
\symidx{\(\operatorname{coker}f\)：模同态 \(f\) 的余核}

这里子模的断言可直接检查：若 \(f(x)=f(y)=0\)，则 \(f(x-y)=0\) 与 \(f(rx)=0\)；若两个元素为 \(f(x),f(y)\)，则其差为 \(f(x-y)\)，\(r\cdot f(x)=f(rx)\)。一个同态是单射，当且仅当核为零，因为 \(f(x)=f(y)\) 等价于 \(x-y\in\ker f\)。它是满射，当且仅当余核为零，因为这恰好表示 \(N=\operatorname{im}f\)。余核把整个像压成零来度量满射性的缺失，并不是从陪域中选出一个与像互补的子模。

\ProseParagraphBreak
记核的包含映射为 \(i_f:\ker f\to M\)。若同态 \(u:P\to M\) 满足 \(fu=0\)，则 \(u(P)\subseteq\ker f\)，故 \(u\) 唯一地分解为 \(u=i_f\bar u\)。余核的商映射满足 \(q_ff=0\)。若同态 \(h:N\to P\) 满足 \(hf=0\)，这正是说 \(h\) 在 \(\operatorname{im}f\) 上为零；由\cref{b2:prop:quotient-module}，它唯一地分解为 \(h=\bar hq_f\)。这两个分解分别由下图的虚线箭头给出：
\[
\begin{tikzcd}[column sep=1.7em,row sep=2em]
P\arrow[r,"u"]\arrow[dr,dashed,"\bar u"']&M\arrow[r,"f"]&N\\
&\ker f\arrow[u,"i_f"']&
\end{tikzcd}
\qquad
\begin{tikzcd}[column sep=1.7em,row sep=2em]
M\arrow[r,"f"]&N\arrow[r,"h"]\arrow[d,"q_f"']&P\\
&\operatorname{coker}f\arrow[ur,dashed,"\bar h"']&
\end{tikzcd}
\]
核与余核在这里都带有指定的映射，分解要求与这些映射相容，而不只是给出一个模的同构类型。

\ProseParagraphBreak
例如同态 \(\mathbb Z\to\mathbb Z\)、\(a\mapsto6a\) 的核为零、像为 \(6\mathbb Z\)、余核为 \(\mathbb Z/6\mathbb Z\)。它是单射，却不是满射。这里 \(6\mathbb Z\) 没有直和补模：任意非零子模 \(L\le\mathbb Z\) 都含有某个非零整数 \(\ell\)，因而 \(6\ell\) 是 \(L\cap6\mathbb Z\) 的非零元素，不能满足直和所要求的交为零。零子模又不能与 \(6\mathbb Z\) 相加得到整个 \(\mathbb Z\)。余核仍然存在，它将全部 \(6\mathbb Z\) 视为零，并不要求从陪域中选出一个补模。

记所有模同态 \(M\to N\) 的集合为 \(\operatorname{Hom}_R(M,N)\)。逐点定义
\[
(f+g)(m)=f(m)+g(m),\qquad (-f)(m)=-f(m).
\]
例如 \((f+g)(rm)=r\cdot f(m)+r\cdot g(m)=r(f+g)(m)\)，所以和仍为模同态；关于加法的其余条件来自 \(N\) 的交换群结构。于是 \(\operatorname{Hom}_R(M,N)\) 总是一个交换群，零元为零映射。
\symidx{\(\operatorname{Hom}_R(M,N)\)：左 \(R\) 模同态组成的交换群}

当定义域与陪域相同，同态还可以复合。记
\[
\operatorname{End}_R(M)\coloneqq\operatorname{Hom}_R(M,M),
\]
以逐点加法和复合乘法 \(fg=f\circ g\) 构成\term[mo-zitongtai-huan]{模自同态环}[endomorphism ring of a module]，单位是恒等映射。复合仍保持加法与 \(R\) 作用，其余环公理与\secref{b2:sec:0201}中的加法自同态环一样逐点核验。

若还要在同态群上定义标量作用，必须检查作用后的映射仍为 \(R\) 线性。新作用可以施加在陪域的值上，也可以先施加在定义域的元素上再取其像；两者都需要与原来的 \(R\) 作用相容。

\begin{proposition}\label{b2:prop:hom-central-action}
设 \(R\) 是环，\(M,N\) 是左 \(R\) 模。环 \(R\) 的中心
\[
Z(R)=\{\,z\in R\mid zr=rz\;\text{对每个}\;r\in R\,\}
\]
通过 \((zf)(m)=z\cdot f(m)\) 使 \(\operatorname{Hom}_R(M,N)\) 成为 \(Z(R)\) 模。特别地，若 \(R\) 交换，同态群以此成为 \(R\) 模。
\end{proposition}
\begin{proof}
中心是含幺交换子环。对 \(z\in Z(R)\)、\(r\in R\)，有
\[
(zf)(rm)=zr\cdot f(m)=rz\cdot f(m)=r\cdot(zf)(m).
\]
因此 \(zf\) 仍为 \(R\) 线性映射；可加性来自 \(N\) 的分配律。对 \(z,z'\in Z(R)\) 逐点计算，\((zz')f=z(z'f)\)、\((z+z')f=zf+z'f\)、\(z(f+g)=zf+zg\)、\(1f=f\)，给出所需模公理。
\end{proof}

不能无条件把 \(Z(R)\) 换成 \(R\)。取 \(M=N={}_RR\)，其中 \(R=M_2(k)\)，令 \(f=\operatorname{id}_R\)，并试图让 \(rf\) 是左乘 \(r\) 的映射。用矩阵单位 \(E_{ij}\) 表示仅第 \((i,j)\) 项为一的矩阵，取 \(r=E_{12}\)、\(s=E_{21}\)，则
\symidx{\(E_{ij}\)：仅指定位置为一的矩阵单位}
\[
(rf)(s)=rs=E_{11},\qquad s(rf)(1_R)=sr=E_{22}.
\]
两者不同，所以 \(rf\) 甚至不是左 \(R\) 模同态。问题在于标量未必能穿过另一个标量，并非加法定义有问题。

\ProseParagraphBreak
设 \(S\) 是环。若陪域 \(N\) 还带左 \(S\) 模结构，且对 \(s\in S\)、\(r\in R\)、\(n\in N\) 满足 \(s(rn)=r(sn)\)，则逐点公式 \((sf)(m)=s\cdot f(m)\) 给 \(\operatorname{Hom}_R(M,N)\) 一个左 \(S\) 模结构。确实，\((sf)(rm)=s(rf(m))=r(sf)(m)\)，所以 \(sf\) 仍为 \(R\) 线性映射；其余模公理由 \(N\) 的作用逐点给出。

\ProseParagraphBreak
另一种左 \(S\) 作用来自定义域。若 \(M\) 同时具有右 \(S\) 模结构，并满足 \((rm)s=r(ms)\)，就称 \(M\) 为 \((R,S)\) \term[shuangmo]{双模}[bimodule]。此时定义
\[
(sf)(m)=f(ms),
\]
这与上一段在 \(N\) 上逐点作用的公式不同。它仍给出左 \(S\) 模结构，因为 \((sf)(rm)=f\bigl((rm)s\bigr)=r\cdot f(ms)\)，而
\[
\bigl(s(tf)\bigr)(m)=(tf)(ms)=f\bigl((ms)t\bigr)=f\bigl(m(st)\bigr)=((st)f)(m).
\]
其余模公理也来自 \(M\) 上的右作用。这里把定义域的右作用变成了同态群上的左作用，次序由 \((ms)t=m(st)\) 保证；它与直接在陪域取值上作用是两种不同的构造。

\begin{proposition}\label{b2:prop:regular-hom-endomorphism}
设 \(R\) 是环，\(M\) 是左 \(R\) 模。取值映射
\[
\operatorname{Hom}_R(R,M)\longrightarrow M,\qquad f\longmapsto f(1_R)
\]
是交换群同构，其逆把 \(m\) 送到 \(r\mapsto rm\)。以复合为乘法的模自同态环满足
\[
\operatorname{End}_R({}_RR)\cong R^{\mathrm{op}}.
\]
\end{proposition}
\begin{proof}
线性条件迫使 \(f(r)=f(r1_R)=r\cdot f(1_R)\)；任意 \(m\) 所给的映射 \(r\mapsto rm\) 确实线性。这两个对应逐点互逆，且保持加法。

当 \(M=R\) 时，\(f\) 因而是某个右乘映射 \(u_a(r)=ra\)。左 \(R\) 线性要求 \(f(sr)=sf(r)\)，参数 \(a=f(1_R)\) 出现在 \(r\) 的右侧；左乘 \(r\mapsto ar\) 则要求 \(asr=sar\)，对一般的 \(a,s\) 并不成立。右乘自同态的复合次序给出
\[
(u_a\circ u_b)(r)=(rb)a=r(ba)=u_{ba}(r).
\]
所以 \(a\mapsto u_a\) 是从反环 \(R^{\mathrm{op}}\) 到该自同态环的保幺环同构。
\end{proof}
\symidx{\(\operatorname{End}_R(M)\)：左 \(R\) 模 \(M\) 的自同态环}

\subsection{模的同构定理}
\label{b2:subsec:0202:03}

\begin{theorem}[第一同构]\label{b2:thm:module-first-isomorphism}
设 \(R\) 是环，\(f:M\to N\) 是左 \(R\) 模之间的同态。它诱导 \(R\) 模同构
\[
M/\ker f\longrightarrow\operatorname{im}f,\qquad m+\ker f\longmapsto f(m).
\]
\end{theorem}
\begin{proof}
把 \(f\) 的陪域改为 \(\operatorname{im}f\)，由\cref{b2:prop:quotient-module}，它在 \(\ker f\) 上为零，因而唯一地下降为所写同态。该同态满射；若 \(f(m)=0\)，则 \(m\in\ker f\)，故源中的陪集为零，所以它也单射。逆映射的线性性由双射同态的一般性质得到。
\end{proof}

这一结果称为\term[mo-diyi-tonggou-dingli]{模的第一同构定理}[first isomorphism theorem for modules]。它所构造的不是 \(M\) 与其像之间任意的同构，而是把全部且仅有的零像元素先模掉。两个向量空间的线性映射满足同一结论，这里不需要维数有限，也不需要标量环是域。

\begin{theorem}[子模对应与商的同构]\label{b2:thm:module-correspondence}
设 \(R\) 是环，\(M\) 是左 \(R\) 模，\(L\le M\)，商映射为 \(q:M\to M/L\)。\(M\) 中包含 \(L\) 的子模 \(U\) 与 \(M/L\) 的子模 \(W\) 通过
\[
U\longmapsto U/L=q(U),\qquad W\longmapsto q^{-1}(W)
\]
一一对应，并保持包含关系。对 \(L\le U\le M\)，有
\[
(M/L)/(U/L)\cong M/U,
\]
同构把 \((m+L)+(U/L)\) 送到 \(m+U\)。
\end{theorem}
\begin{proof}
同态的像与逆像保持子模条件，且 \(q^{-1}(W)\) 总包含 \(L\)。若 \(U\supseteq L\)，则 \(q(m)\in q(U)\) 等价于某个 \(u\in U\) 满足 \(m-u\in L\)，也就等价于 \(m\in U\)。故 \(q^{-1}(q(U))=U\)。由于 \(q\) 满射，对任意 \(W\) 有 \(q(q^{-1}(W))=W\)。两种操作保持包含关系，得到所述对应。

映射 \(M\to M/U\)、\(m\mapsto m+U\) 在 \(L\) 上为零。商映射的泛性质使它下降为满同态 \(M/L\to M/U\)，其核是所有满足 \(m\in U\) 的陪集，即 \(U/L\)。应用\cref{b2:thm:module-first-isomorphism}就得到第二个同构及其元素公式。
\end{proof}

\begin{theorem}[第二同构]\label{b2:thm:module-second-isomorphism}
设 \(R\) 是环，\(M\) 是左 \(R\) 模。若 \(U,V\le M\)，则 \(U+V=\{\,u+v\mid u\in U,\ v\in V\,\}\) 与 \(U\cap V\) 均为子模，并有 \(R\) 模同构
\[
U/(U\cap V)\longrightarrow(U+V)/V,\qquad
u+(U\cap V)\longmapsto u+V.
\]
\end{theorem}
\begin{proof}
和中的两个元素相减后仍为来自 \(U,V\) 的元素之和，标量乘法也保持这种形式；交中的元素同时满足两个子模的封闭条件。因此它们都是子模。考虑同态 \(U\to(U+V)/V\)、\(u\mapsto u+V\)。每个陪集 \((u+v)+V=u+V\)，故它满射；零像的条件为 \(u\in V\)，故核为 \(U\cap V\)。再用\cref{b2:thm:module-first-isomorphism}即可。
\end{proof}

后两个结果分别包含\term[mo-disan-tonggou-dingli]{模的第三同构定理}[third isomorphism theorem for modules]与\term[mo-dier-tonggou-dingli]{模的第二同构定理}[second isomorphism theorem for modules]。两者都先构造一个自然满同态，再计算它的核，最后应用第一同构定理：第二同构使用 \(U\to(U+V)/V\)，第三同构使用 \(M/L\to M/U\)。商模公式记录的是这两个具体映射的结果。作为一个算例，取 \(M=\mathbb Z\)、\(U=4\mathbb Z\)、\(V=6\mathbb Z\)。此时 \(U+V=2\mathbb Z\)、\(U\cap V=12\mathbb Z\)，得到
\[
4\mathbb Z/12\mathbb Z\cong2\mathbb Z/6\mathbb Z.
\]
同构把 \(4+12\mathbb Z\) 送到 \(4+6\mathbb Z\)；两边都有三个元素，但对应公式还告诉我们具体的生成元怎样相配。

% !TeX root = ../../Book2.tex
% 纲要标题：### 2.3 生成、零化子与挠性

\section{生成、零化子与挠性}
\label{b2:sec:0203}

商模的运算已经在\cref{b2:prop:quotient-module}中建立。现在转而考察怎样指定子模，以及模掉一个子模究竟加入了哪些关系。与此相连的另一个问题是：哪些标量在某个元素上，或在整个模上作用为零？这两种零化条件不同，必须分别记录。

% 纲要要点（供目录核对）：
% * 子模生成与关系
% * 零化子与忠实模；正文证明商左模的零化子为左理想中的最大双边理想
% * 模上的零因子和挠

\subsection{子模生成与关系}
\label{b2:subsec:0203:01}

\begin{proposition}\label{b2:prop:generated-submodule}
设 \(R\) 是环，\(M\) 是左 \(R\) 模。任意一族 \(M\) 的子模的交仍为子模，空族的交约定为 \(M\)。对任意子集 \(S\subseteq M\)，存在包含 \(S\) 的最小子模，称为 \(S\) \term[shengcheng-zimo]{生成的子模}[submodule generated by a subset]，记为 \(\langle S\rangle_R\)。具体地，
\[
\langle S\rangle_R=
\left\{\,\sum_{j=1}^n r_js_j\ \middle|\ n\ge0,\ r_j\in R,\ s_j\in S\,\right\},
\]
其中空和为零。当 \(S=\{m\}\) 时记为 \(Rm\)；有 \(M=Rm\) 的模称为\term[xunhuan-mo]{循环模}[cyclic module]。
\end{proposition}
\begin{proof}
属于每个子模的元素，其差与任意标量倍仍属于每个子模，所以交为子模。所有包含 \(S\) 的子模形成非空族，因为 \(M\) 在其中；取交便得到最小子模。

右侧的有限和集合包含零及 \(S\)。两个有限和相减仍是有限和，左乘标量 \(a\) 后各系数变为 \(ar_j\)，也仍有同样形式，因此该集合是子模。反过来，任何包含 \(S\) 的子模都包含所有这些有限线性组合，所以它正是最小子模。
\end{proof}
\symidx{\(\langle S\rangle_R\)、\(Rm\)：集合或单个元素生成的子模}

对一族子模 \((L_i)_{i\in I}\)，其\term[zimo-de-he]{和}[sum of submodules]是并集生成的子模，记为 \(\sum_{i\in I}L_i\)。它的每个元素都是有限多个 \(L_i\) 中的元素之和；允许指标集无限，不等于允许在模中作无限加法。例如在 \(\mathbb Z\) 模中，\(4\mathbb Z+6\mathbb Z=2\mathbb Z\)，因为所有和都是偶数，而 \(2=6-4\)。一般并集未必是子模：\(4\) 与 \(6\) 属于 \(4\mathbb Z\cup6\mathbb Z\)，它们的和 \(10\) 却不属于其中。

\begin{definition}\label{b2:def:finitely-generated-module}
设 \(R\) 是环，\(M\) 是左 \(R\) 模。若存在有限个元素 \(m_1,\ldots,m_n\in M\)，使 \(M=Rm_1+\cdots+Rm_n\)，就称 \(M\) 为\term[youxian-shengcheng-mo]{有限生成模}[finitely generated module]，并称这些元素为一组生成元。\footnote{“有限模”是“有限生成模”的简称；这里的“有限”不表示模的底集合有限。}零模允许由空集生成。
\end{definition}

\(\mathbb Z\) 作为自身上的模由 \(1\) 生成，底集合却无限；\(\mathbb Q\) 作为 \(\mathbb Q\) 模由 \(1\) 生成，作为 \(\mathbb Z\) 模却不有限生成。后一断言可以直接验证：若有限多个有理数的分母都整除正整数 \(d\)，它们的整数线性组合都在 \(d^{-1}\mathbb Z\) 中，而 \(1/(2d)\) 不在其中。因此是否有限生成总是相对于指定标量环而言。

\ProseParagraphBreak
给出生成元并没有给出唯一坐标。称一个等式 \(r_1m_1+\cdots+r_nm_n=0\) 为这组生成元之间的一个关系。若 \(L\le M\)，则在 \(M/L\) 中一个线性组合为零，恰好表示它在 \(M\) 中属于 \(L\)。例如，在
\[
M=\mathbb Z^2/\mathbb Z(2,3)
\]
中，标准向量的像 \(\bar e_1,\bar e_2\) 满足 \(2\bar e_1+3\bar e_2=0\)。商映射 \(q:\mathbb Z^2\to M\) 把 \((a,b)\) 送到 \(a\bar e_1+b\bar e_2\)，故
\[
a\bar e_1+b\bar e_2=0
\quad\Longleftrightarrow\quad
(a,b)\in\ker q=\mathbb Z(2,3).
\]
因此每个关系的系数对都为 \(n(2,3)\)，其中 \(n\in\mathbb Z\)，即全部关系都由 \(2\bar e_1+3\bar e_2=0\) 这一条生成。这里 \(\mathbb Z^2\) 采用逐坐标运算。映射 \((a,b)\mapsto3a-2b\) 满射到 \(\mathbb Z\)，因为 \((1,1)\) 映到 \(1\)；它的核由 \(3a=2b\) 给出，即 \(\mathbb Z(2,3)\)。由\cref{b2:thm:module-first-isomorphism}，\(M\cong\mathbb Z\)。

\begin{proposition}\label{b2:prop:cyclic-module-quotient}
设 \(R\) 是环，\(M\) 是左 \(R\) 模，\(m\in M\)。满同态 \(R\to Rm\)、\(r\mapsto rm\) 的核是左理想
\[
I_m=\{\,r\in R\mid rm=0\,\},
\]
并诱导左 \(R\) 模同构 \(R/I_m\cong Rm\)。所以每个循环左模均是左正则模的一个商。
\end{proposition}
\begin{proof}
映射可加，且把 \(ar\) 送到 \((ar)m=a(rm)\)，故为左模同态，像按定义为 \(Rm\)。核是 \(I_m\)，作为正则模的子模它是左理想。再由\cref{b2:thm:module-first-isomorphism}得到同构。
\end{proof}

这里的 \(R/I_m\) 是商模；只有 \(I_m\) 为双边理想时，陪集乘法才给出商环。通过一个元素描述循环模，不需要把它先做成一个环。

\subsection{零化子与忠实模}
\label{b2:subsec:0203:02}

\begin{definition}\label{b2:def:annihilator-faithful}
设 \(R\) 是环，\(M\) 是左 \(R\) 模，\(S\subseteq M\)。定义 \(S\) 的\term[linghuazi]{零化子}[annihilator]为
\[
\operatorname{Ann}_R(S)=\{\,r\in R\mid rs=0\;\text{对每个}\;s\in S\,\}.
\]
当 \(S=\{m\}\) 时记为 \(\operatorname{Ann}_R(m)\)。若 \(\operatorname{Ann}_R(M)=0\)，就称 \(M\) 为\term[zhongshi-mo]{忠实模}[faithful module]。
\end{definition}
\symidx{\(\operatorname{Ann}_R(S)\)：子集 \(S\) 在环 \(R\) 中的零化子}

单个元素的零化子正是\cref{b2:prop:cyclic-module-quotient}中的核 \(I_m\)，所以该命题也可写为
\[
I_m=\operatorname{Ann}_R(m),\qquad
Rm\cong R/\operatorname{Ann}_R(m).
\]
这里仍是左模同构；在非交换环上，\(\operatorname{Ann}_R(m)\) 只保证是左理想。

\ProseParagraphBreak
\cref{b2:def:vector-annihilator}中的 \(U^\circ\) 位于对偶空间，收集在 \(U\) 上为零的线性泛函；这里的 \(\operatorname{Ann}_R(S)\) 位于标量环，收集在 \(S\) 上作用为零的标量。名称相同，但所收集的对象及所在空间不同。

\begin{proposition}\label{b2:prop:annihilator-ideals}
设 \(R\) 是含幺环，\(M\) 是幺性左 \(R\) 模。对任意子集 \(S\subseteq M\)，\(\operatorname{Ann}_R(S)\) 是左理想。如果 \(S\) 本身为子模，则它是双边理想。特别地，\(\operatorname{Ann}_R(M)\) 是作用表示 \(\rho:R\to\operatorname{End}_{\mathbb Z}(M)\)、\(\rho(r)(m)=rm\) 的核，因而 \(M\) 忠实当且仅当 \(\rho\) 单射。

若 \(I\) 是 \(R\) 的双边理想，原有 \(R\) 作用能通过商环 \(R/I\) 定义为
\[
(r+I)m=rm
\]
当且仅当 \(I\subseteq\operatorname{Ann}_R(M)\)。成立时这个 \(R/I\) 模结构唯一。

另设 \(L\) 为 \(R\) 的左理想，则商左模的零化子为
\[
\operatorname{Ann}_R(R/L)=\{\,r\in R\mid rR\subseteq L\,\},
\]
它是包含在 \(L\) 中的最大双边理想。
\end{proposition}
\begin{proof}
若 \(a,b\) 零化每个 \(s\in S\)，则 \((a-b)s=0\)；任意 \(r\in R\) 满足 \((ra)s=r(as)=0\)，故零化子为左理想。若 \(S\) 是子模，则 \(rs\in S\)，因而 \((ar)s=a(rs)=0\)，又得到右乘封闭性。由\cref{b2:prop:module-representation}的表示公式，\(\rho(a)=0\) 恰好表示 \(aM=0\)，故其核等于全模零化子。

若商环作用有上述形式，取 \(a\in I\)，便有 \(am=(a+I)m=0\)。反过来，若 \(IM=0\)，当 \(r-r'\in I\) 时 \(rm-r'm=0\)，所以作用不依赖代表元。商环的加法、乘法和单位元运算使四条模公理分别从原来的作用下降。由于每个商环标量都是某个 \(r+I\)，规定相容性后没有别的选择。

标量 \(r\) 零化 \(R/L\) 的每个陪集 \(a+L\)，当且仅当对全部 \(a\in R\) 有 \(ra\in L\)，即 \(rR\subseteq L\)。取 \(a=1_R\) 可知这个零化子包含在 \(L\) 中，而全模零化子已证明为双边理想。若 \(J\subseteq L\) 为双边理想，则每个 \(r\in J\) 都满足 \(rR\subseteq J\subseteq L\)，故 \(J\subseteq\operatorname{Ann}_R(R/L)\)，证明其最大性。
\end{proof}

忠实性测量的是标量能否由它们在模上的作用区分。由 \(\ker\rho=0\)，它等价于
\[
r\ne s\quad\Longrightarrow\quad
\exists m\in M,\qquad rm\ne sm.
\]
事实上，\(rm=sm\) 对所有 \(m\) 成立，恰好表示 \(r-s\in\operatorname{Ann}_R(M)\)。每个模都可因此看成 \(R/\operatorname{Ann}_R(M)\) 上的忠实模：若一个陪集零化所有元素，其代表元已属于全模零化子。这个商环只除去对整个模完全不起作用的标量。正则模 \({}_RR\) 忠实，因为零化 \(1_R\) 的标量只能为零；零模的零化子是 \(R\)，所以零模只在 \(R\) 为零环时忠实。

\ProseParagraphBreak
单个元素的零化子可能更大，且未必是双边理想。令 \(R=M_2(k)\) 作用在列向量 \(M=k^2\) 上。一个矩阵零化所有列向量，意味着它的两列都是零，故 \(M\) 忠实。但对 \(e_1=(1,0)^{\mathsf T}\)，
\[
\operatorname{Ann}_R(e_1)
=\left\{\,\begin{pmatrix}0&a\\0&b\end{pmatrix}\ \middle|\ a,b\in k\,\right\}.
\]
它包含 \(E_{12}\)，却不包含 \(E_{12}E_{21}=E_{11}\)，所以不是右理想。忠实性说的是没有一个非零标量零化全部元素，并不要求每个非零元素都具有零零化子。

\subsection{模上的零因子和挠}
\label{b2:subsec:0203:03}

先设 \(R\) 为交换环。给定模 \(M\)，一个标量 \(r\) 是否可以消去，要看映射 \(M\to M\)、\(m\mapsto rm\) 是否单射；这不仅取决于 \(R\)，也取决于 \(M\)。

\begin{definition}\label{b2:def:module-zero-divisor}
设 \(R\) 为交换环，\(M\) 为 \(R\) 模，\(r\in R\)。若存在 \(0\ne m\in M\) 使 \(rm=0\)，就称 \(r\) 为 \(M\) 上的\term[mo-shang-lingyinzi]{零因子}[zero divisor on a module]。本书在这个集合中包括满足条件的零标量：当 \(M\ne0\) 时，\(0\) 是 \(M\) 上的零因子。若 \(r\) 不是 \(M\) 上的零因子，就称 \(r\) 为 \(M\) 上的\term[mo-shang-zhengze-yuan]{正则元}[regular element on a module]。
\end{definition}

对交换环 \(R\) 的非零元素 \(r\)，第一卷中环零因子的定义正好是取正则模 \({}_RR\) 后的模零因子条件：存在 \(0\ne a\in R\) 使 \(ra=0\)。零标量的约定则有所不同：第一卷的环零因子不包括 \(0\)，而本节的模零因子在模非零时包括 \(0\)。例如，\(2\) 不是环 \(\mathbb Z\) 上的零因子，却是 \(\mathbb Z/6\mathbb Z\) 上的零因子，因为 \(2\bar3=0\)。在这个模上，\(2\) 与 \(3\) 都是零因子，但它们的和 \(5\) 可逆地作用，故模上零因子的集合一般不是理想。零模上没有零因子，因为它没有非零元素。

\begin{definition}\label{b2:def:torsion-module}
设 \(R\) 是整环，即非零交换环且没有非零零因子，\(M\) 是 \(R\) 模。若 \(m\in M\) 被某个 \(0\ne r\in R\) 零化，就称 \(m\) 为\term[naoyuan]{挠元}[torsion element]。所有挠元的集合记为 \(t_R(M)\)。若 \(t_R(M)=M\)，就称 \(M\) 为\term[naomo]{挠模}[torsion module]；若 \(t_R(M)=0\)，就称 \(M\) 为\term[wunao-mo]{无挠模}[torsion-free module]。
\end{definition}
\symidx{\(t_R(M)\)：整环 \(R\) 上模 \(M\) 的挠子模}

\begin{proposition}\label{b2:prop:torsion-submodule}
设 \(R\) 是整环，\(M,N\) 是 \(R\) 模，则 \(t_R(M)\) 是 \(M\) 的子模，称为挠子模，且 \(M/t_R(M)\) 无挠。每个 \(R\) 模同态 \(f:M\to N\) 都把 \(t_R(M)\) 映入 \(t_R(N)\)。
\end{proposition}
\begin{proof}
零元由 \(1_R\ne0\) 零化。若 \(ax=0\)、\(by=0\)，其中 \(a,b\ne0\)，则整环条件保证 \(ab\ne0\)，而交换性给出
\[
ab(x-y)=b(ax)-a(by)=0.
\]
故挠元对差封闭。任意 \(r\in R\) 满足 \(a(rx)=r(ax)=0\)，所以也对标量作用封闭。

设 \(0\ne a\in R\)，且 \(a(m+t_R(M))=0\)。这表示 \(am\in t_R(M)\)，因而某个 \(0\ne b\in R\) 满足 \(bam=0\)。由于 \(ba\ne0\)，\(m\) 本身是挠元，其陪集为零，证明商模无挠。最后，若 \(am=0\)，则 \(a\cdot f(m)=f(am)=0\)，同一个非零标量也零化 \(f(m)\)。
\end{proof}

整环假设在证明中有明确用途：两个非零零化标量相乘后仍须非零。若在带零因子的环上仍用“被某个非零标量零化”定义挠元，它们未必组成子模。取 \(R=M=\mathbb Z/6\mathbb Z\)，\(\bar2\) 与 \(\bar3\) 分别被非零标量 \(\bar3\) 与 \(\bar2\) 零化，但和 \(\bar5\) 是单位，不被任何非零标量零化。本章的挠模结论因此限于整环；在一般环上还可以采用正则标量定义另一种挠条件，使用时须另作约定。

\ProseParagraphBreak
挠性与忠实性也不能混同。在整环 \(R\) 上，\(M\) 为挠模的条件是
\[
\forall m\in M,\quad
\exists\,0\ne r_m\in R,\qquad r_m m=0;
\]
而 \(M\) 非忠实的条件是
\[
\exists\,0\ne r\in R,\quad
\forall m\in M,\qquad rm=0.
\]
前者允许零化标量随元素变化，后者要求同一个非零标量零化整个模。\(\mathbb Q/\mathbb Z\) 作为 \(\mathbb Z\) 模是挠模，因为 \(a/b+\mathbb Z\) 被非零整数 \(b\) 零化；它却是忠实的。事实上，对任意非零整数 \(n\)，取 \(q=1/(|n|+1)\)，则 \(nq\notin\mathbb Z\)，故 \(n\) 不能零化整个模。

\ProseParagraphBreak
有限生成性使这两种量词条件一致。若整环上的挠模由 \(m_1,\ldots,m_n\) 生成，分别选 \(0\ne a_i\) 使 \(a_im_i=0\)，则整环条件保证乘积 \(a=a_1\cdots a_n\) 非零。这个乘积零化每个生成元，由交换性也零化它们的所有线性组合，即 \(aM=0\)。零模允许空生成集，此时取空乘积 \(a=1_R\ne0\)，结论仍成立。反过来，任何非忠实模的统一非零零化标量也零化每个元素，因此该模为挠模。于是整环上的有限生成模为挠模，当且仅当它非忠实。\(\mathbb Q/\mathbb Z\) 的例子说明，这里不能略去有限生成条件。

% !TeX root = ../../Book2.tex
% 纲要标题：### 2.4 直和与直积

\section{直和与直积}
\label{b2:sec:0204}

把几个模放在一起，最直接的办法是组成有序元组，再逐坐标运算。有限多个因子时，这同时解决了两类问题：给出映入每个坐标的同态，便能得到映入整个元组模的同态；给出每个坐标模映到另一个模的同态，也能把像相加。当因子有无限多个时，第二种做法遇到无限求和的问题。直积与直和的区分就来自这两种映射要求的不同。

% 纲要要点（供目录核对）：
% * 任意族的直和及直积
% * 有限情形的一致与无限情形的差异
% * 同态进入直积、离开直和的泛性质
% * 有限生成模映入直和时的共同有限支集与同态群同构
% ---

\subsection{任意族的直和与直积}
\label{b2:subsec:0204:01}

\begin{definition}\label{b2:def:module-product-sum}
设 \(R\) 是环，\(I\) 是集合，\((M_i)_{i\in I}\) 是一族左 \(R\) 模。它们的\term[mo-zhiji]{直积}[direct product of modules]为
\[
\prod_{i\in I}M_i=\{\,(m_i)_{i\in I}\mid m_i\in M_i\text{ 对每个 }i\in I\,\},
\]
加法与标量乘法均逐坐标定义。元素 \(m=(m_i)\) 的\term[zhiji]{支集}[support]为 \(\operatorname{supp}(m)=\{\,i\in I\mid m_i\ne0\,\}\)。有限支集元素组成的子模
\[
\bigoplus_{i\in I}M_i
=\left\{\,m\in\prod_{i\in I}M_i\ \middle|\ \operatorname{supp}(m)\text{ 有限}\,\right\}
\]
称为这族模的\term[mo-zhihe]{直和}[direct sum of modules]。当所有因子均为 \(M\) 时，分别简记为 \(M^I\) 与 \(M^{(I)}\)。
\end{definition}
\symidx{\(\operatorname{supp}(m)\)：模直积元素的非零坐标支集}
\symidx{\(\prod_{i\in I}M_i\)、\(M^I\)：模族的直积}
\symidx{\(\bigoplus_{i\in I}M_i\)、\(M^{(I)}\)：模族的直和}

直积中全部模公理可以逐坐标检查。例如，\(r(sm_i)=(rs)m_i\) 对每个 \(i\) 成立，就表示两个相应元组相等。对直和还要检查有限支集条件：\(x-y\) 的支集包含在 \(\operatorname{supp}(x)\cup\operatorname{supp}(y)\) 中，\(rx\) 的支集包含在 \(\operatorname{supp}(x)\) 中，二者均有限，所以它是子模。空指标集只有一个空元组，其直积与直和都为零模。

\ProseParagraphBreak
记坐标投影为 \(p_i:\prod_jM_j\to M_i\)。直和有坐标嵌入 \(\iota_i:M_i\to\bigoplus_jM_j\)，把元素放在第 \(i\) 坐标，其余坐标取零。二者都是模同态，而且
\[
p_i\iota_j=
\begin{cases}
\operatorname{id}_{M_i},&i=j,\\
0,&i\ne j,
\end{cases}
\]
其中投影在直和上取限制。每个直和元素都有唯一的有限展开
\[
m=\sum_{i\in\operatorname{supp}(m)}\iota_i(m_i).
\]
这项有限性将使离开直和的映射可以通过逐项求和定义。

\subsection{有限与无限情形的比较}
\label{b2:subsec:0204:02}

当 \(I\) 有限时，所有元组都有有限支集，因而上述直和与直积就是同一个模。特别地，两因子情形可写 \(M\oplus N=M\times N\)。当 \(I=\mathbb N\) 且各因子为非零环 \(R\) 的正则模时，直积中的元素 \((1,1,1,\ldots)\) 不在直和中。故自然包含
\[
R^{(\mathbb N)}\longrightarrow R^{\mathbb N}
\]
不是满射。这里比较的是直和与直积之间给定的自然映射。

\ProseParagraphBreak
还须区别上述由外部元组构造的直和与一个模内部的分解。若已给 \(M\) 的子模 \(L_i\)，可以考虑它们的和是否包含重叠的表示。

\begin{proposition}\label{b2:prop:internal-direct-sum}
设 \(R\) 是环，\(M\) 是左 \(R\) 模，\((L_i)_{i\in I}\) 是 \(M\) 的一族子模。求和映射
\[
\sigma:\bigoplus_{i\in I}L_i\longrightarrow M,\qquad
(\ell_i)\longmapsto\sum_i\ell_i
\]
是同态，像为 \(\sum_iL_i\)。它是同构，当且仅当 \(M=\sum_iL_i\)，且对每个 \(i\in I\) 都有
\[
L_i\cap\sum_{j\ne i}L_j=0.
\]
成立时，每个 \(m\in M\) 都有唯一的有限表达 \(m=\sum_i\ell_i\)、\(\ell_i\in L_i\)，此时称 \(M\) 为这些子模的\term[nei-zhihe]{内直和}[internal direct sum]，也记 \(M=\bigoplus_iL_i\)。
\end{proposition}
\begin{proof}
直和中的有限支集使求和有意义，有限和的分配律给出 \(\sigma\) 的线性性，像由子模和的定义确定。假设交条件成立，且 \(\sum_i\ell_i=0\)。对于每个非零坐标 \(i\)，有 \(\ell_i=-\sum_{j\ne i}\ell_j\)，属于题设的交，因此为零。于是核为零；再有 \(M=\sum_iL_i\) 就得到满射和同构。

反过来，若 \(\sigma\) 单射，而 \(x\in L_i\cap\sum_{j\ne i}L_j\)，写 \(x=\sum_{j\ne i}\ell_j\)，就得到一个有限支集元组，其第 \(i\) 项为 \(x\)，其余有关项为 \(-\ell_j\)，且像为零。单射性迫使 \(x=0\)。最后，同构使每个 \(m\) 对应唯一元组，这正是有限表达的唯一性。
\end{proof}

两因子时条件简化为 \(M=U+V\)、\(U\cap V=0\)。三个或更多因子时，仅要求两两交为零还不够：\(k^2\) 中三条直线 \(ke_1\)、\(ke_2\)、\(k(e_1+e_2)\) 两两交为零，却有来自三个子空间的非平凡关系 \(e_1+e_2-(e_1+e_2)=0\)。一般模也未必像向量空间那样总能找到补子模。例如 \(2\mathbb Z\le\mathbb Z\) 没有直和补模。若 \(\mathbb Z=2\mathbb Z\oplus L\)，把商映射 \(q:\mathbb Z\to\mathbb Z/2\mathbb Z\) 限制在 \(L\) 上：由 \(L\cap2\mathbb Z=0\)，其核为零；由 \(\mathbb Z=2\mathbb Z+L\)，每个陪集都有 \(L\) 中的代表元，所以限制映射满射。因此 \(L\cong\mathbb Z/2\mathbb Z\)。但 \(\mathbb Z\) 的非零元素都没有有限加法阶，不能含有这样的子模。

\begin{example}
在 \(R=\mathbb Z\) 时，直和与直积甚至不可能有抽象群同构。\(\mathbb Z^{(\mathbb N)}\) 是由前 \(n\) 个坐标支撑的可数集合（各与 \(\mathbb Z^n\) 对应）的可数并，故仍可数；而 \(\mathbb Z^{\mathbb N}\) 包含不可数的零一序列集合 \(\{0,1\}^{\mathbb N}\)。后者的不可数性由对角法得到：任意列举的第 \(n\) 个序列若在第 \(n\) 位取 \(a_n\)，就令新序列在该位取 \(1-a_n\)。
\end{example}

\subsection{直积与直和的泛性质}
\label{b2:subsec:0204:03}

映入直积的同态由各个坐标投影确定，映出直和的同态由它在各个坐标嵌入上的限制确定；下面两个泛性质分别表达这两种方向。

\begin{theorem}\label{b2:thm:module-product-universal}
设 \(R\) 是环，\((M_i)_{i\in I}\) 是一族左 \(R\) 模，\(N\) 也是左 \(R\) 模。对每个 \(i\in I\) 给定同态 \(f_i:N\to M_i\)，并记 \(p_i:\prod_{j\in I}M_j\to M_i\) 为坐标投影。此时存在唯一同态
\[
f:N\longrightarrow\prod_{i\in I}M_i
\]
满足 \(p_if=f_i\) 对每个 \(i\) 成立。它由 \(f(n)=(f_i(n))_{i\in I}\) 给出。
等价地，对每个 \(i\in I\)，下图交换：
\[
\begin{tikzcd}[column sep=3.5em,row sep=2.5em]
N\arrow[r,dashed,"f"]\arrow[dr,"f_i"']&\prod_{j\in I}M_j\arrow[d,"p_i"]\\
&M_i
\end{tikzcd}
\]
\end{theorem}
\begin{proof}
坐标 \(f_i(n)\) 属于 \(M_i\)，所以公式定义了直积中的元素。每个坐标上的加法与标量等式都由 \(f_i\) 的线性性保证，故 \(f\) 线性。条件 \(p_if=f_i\) 已经逐个确定了 \(f(n)\) 的坐标，因此没有另一种选择。
\end{proof}

\begin{theorem}\label{b2:thm:module-sum-universal}
设 \(R\) 是环，\((M_i)_{i\in I}\) 是一族左 \(R\) 模，\(N\) 也是左 \(R\) 模。对每个 \(i\in I\) 给定同态 \(g_i:M_i\to N\)，并记 \(\iota_i:M_i\to\bigoplus_{j\in I}M_j\) 为坐标嵌入。此时存在唯一同态
\[
g:\bigoplus_{i\in I}M_i\longrightarrow N
\]
满足 \(g\iota_i=g_i\) 对每个 \(i\) 成立。它由有限和
\[
g((m_i))=\sum_{i\in\operatorname{supp}(m)}g_i(m_i)
\]
给出。
等价地，对每个 \(i\in I\)，下图交换：
\[
\begin{tikzcd}[column sep=3.5em,row sep=2.5em]
M_i\arrow[r,"\iota_i"]\arrow[dr,"g_i"']&\bigoplus_{j\in I}M_j\arrow[d,dashed,"g"]\\
&N
\end{tikzcd}
\]
\end{theorem}
\begin{proof}
每个元组的非零坐标只有有限个，故和有意义，添加零坐标不改变值。检验加法时，可以把两个支集同时扩充到它们的有限并集；对这个共同有限集逐项使用 \(g_i(m_i+n_i)=g_i(m_i)+g_i(n_i)\)，就得到可加性。对标量乘法也可保留原支集并加入可能出现的零项，由 \(g_i(rm_i)=r\cdot g_i(m_i)\) 得到线性性。公式直接满足 \(g\iota_i=g_i\)。由于直和中每个元素是 \(\iota_i(m_i)\) 的有限和，这些限制也唯一地决定 \(g\)。
\end{proof}

用同态群表示，这两个泛性质给出交换群同构
\[
\operatorname{Hom}_R\!\left(N,\prod_iM_i\right)
\cong\prod_i\operatorname{Hom}_R(N,M_i),
\]
\[
\operatorname{Hom}_R\!\left(\bigoplus_iM_i,N\right)
\cong\prod_i\operatorname{Hom}_R(M_i,N).
\]
右侧都是直积，因为允许任意一族同态，并不要求其中仅有有限多个非零。

\ProseParagraphBreak
泛性质还说明带有这些映射的对象在唯一同构意义下确定。以直积为例，若 \(P\) 连同同态 \(q_i:P\to M_i\) 也满足\cref{b2:thm:module-product-universal}中的要求，两边的泛性质分别给出 \(u:P\to\prod_iM_i\) 与 \(v:\prod_iM_i\to P\)，满足 \(p_iu=q_i\)、\(q_iv=p_i\)。对每个 \(i\)，这两条等式表示下图的两条三角形路径都交换：
\[
\begin{tikzcd}[column sep=3em,row sep=2.5em]
P\arrow[rr,bend left=15,"u"]\arrow[dr,"q_i"']&&
\prod_jM_j\arrow[ll,bend left=15,"v"]\arrow[dl,"p_i"]\\
&M_i&
\end{tikzcd}
\]
于是 \(q_i(vu)=p_iu=q_i\)，而 \(p_i(uv)=q_iv=p_i\)。因此 \(vu\) 与 \(\operatorname{id}_P\) 有相同的所有 \(q_i\) 分量，\(uv\) 与 \(\operatorname{id}_{\prod_iM_i}\) 有相同的所有 \(p_i\) 分量；各自泛性质中的唯一性给出 \(vu=\operatorname{id}_P\)、\(uv=\operatorname{id}_{\prod_iM_i}\)。这也证明 \(u\) 是同构，而满足 \(p_iu=q_i\) 的映射原本就唯一，所以相容的同构唯一。

\ProseParagraphBreak
直和的情形把投影条件换成嵌入条件。若 \(Q\) 连同 \(j_i:M_i\to Q\) 也满足\cref{b2:thm:module-sum-universal}，两边的泛性质给出 \(a:\bigoplus_iM_i\to Q\)、\(b:Q\to\bigoplus_iM_i\)，满足 \(a\iota_i=j_i\)、\(bj_i=\iota_i\)。于是 \((ba)\iota_i=\iota_i\)、\((ab)j_i=j_i\)，唯一性分别给出 \(ba=\operatorname{id}_{\bigoplus_iM_i}\)、\(ab=\operatorname{id}_Q\)。因此满足坐标嵌入相容条件的同构也唯一。

\ProseParagraphBreak
有限多个因子时，直和与直积相同，所以两个方向的泛性质都适用。无限情形中则不能交换方向而不加条件。例如，对 \(P=\mathbb Z^{\mathbb N}\)、\(S=\mathbb Z^{(\mathbb N)}\)，商映射 \(q:P\to P/S\) 在每个单坐标嵌入的像上都为零，却不是零映射，因为 \((1,1,1,\ldots)\notin S\)。因此，映出无限直积的同态未必由它在各个单坐标上的值决定。

\ProseParagraphBreak
另一方面，一族 \(f_i:N\to M_i\) 所确定的直积映射能够落入直和，当且仅当
\[
\forall n\in N,\qquad\{\,i\in I\mid f_i(n)\ne0\,\}\;\text{有限}.
\]
这不要求整族同态中只有有限多个非零。例如，取非零环 \(R\)，令 \(N=R^{(\mathbb N)}\)，并令 \(f_i:N\to R\) 为第 \(i\) 个坐标投影。每个 \(n\) 的支集有限，所以只有有限多个 \(f_i(n)\ne0\)；但 \(f_i\) 把第 \(i\) 个坐标为 \(1\)、其余为零的元组映为 \(1\)，故每个 \(f_i\) 都非零。有限生成性使这两个有限性条件一致。

\begin{proposition}\label{b2:prop:finite-hom-direct-sum}
设 \(R\) 是含幺环，\(M\) 是有限生成左 \(R\) 模，\((N_i)_{i\in I}\) 是任意一族左 \(R\) 模。每个同态 \(f:M\to\bigoplus_{i\in I}N_i\) 都存在有限集 \(F\subseteq I\)，使
\[
f(M)\subseteq\bigoplus_{i\in F}N_i,
\]
其中右侧通过坐标嵌入看作整个直和的子模。记 \(\iota_i:N_i\to\bigoplus_{j\in I}N_j\) 为坐标嵌入，则映射
\[
\begin{aligned}
\Phi:\bigoplus_{i\in I}\operatorname{Hom}_R(M,N_i)
&\longrightarrow\operatorname{Hom}_R\!\left(M,\bigoplus_{i\in I}N_i\right),\\
(f_i)&\longmapsto\left(m\longmapsto\sum_i\iota_i(f_i(m))\right)
\end{aligned}
\]
是自然的交换群同构。
\end{proposition}
\begin{proof}
取有限生成元 \(m_1,\ldots,m_s\)，令
\[
F=\bigcup_{j=1}^s\operatorname{supp}(f(m_j)).
\]
每个支集有限，所以 \(F\) 有限。若 \(m=\sum_{j=1}^sr_jm_j\)，则 \(f(m)=\sum_{j=1}^sr_j\cdot f(m_j)\) 在 \(F\) 之外的坐标都为零，因为逐坐标标量乘法与加法不会产生新的非零坐标。因此 \(f(M)\subseteq\bigoplus_{i\in F}N_i\)。若 \(M=0\)，可以取空生成集，此时 \(F=\varnothing\)，结论仍成立。

定义 \(\Phi\) 时，\((f_i)\) 只有有限多个非零分量，因此式中的求和只涉及同一个有限集，得到的映射是 \(R\)-线性的。逐项相加也表明 \(\Phi\) 是交换群同态。记 \(p_i\) 为整个直和上的第 \(i\) 个坐标投影；由 \(p_i\iota_j=0\)（\(i\ne j\)）及 \(p_i\iota_i=\operatorname{id}_{N_i}\)，有 \(p_i\Phi((f_j))=f_i\)。所以 \(\Phi((f_i))=0\) 迫使每个 \(f_i=0\)，即 \(\Phi\) 单射。

对任意 \(f:M\to\bigoplus_iN_i\)，第一部分给出的有限集 \(F\) 满足 \(p_if=0\) 对所有 \(i\notin F\) 成立。因此 \((p_if)_{i\in I}\) 属于 \(\bigoplus_i\operatorname{Hom}_R(M,N_i)\)。每个 \(f(m)\) 的有限坐标展开又给出
\[
\Phi((p_if)_{i\in I})(m)=\sum_i\iota_i(p_if(m))=f(m),
\]
故 \(\Phi\) 满射，逆映射就是 \(f\mapsto(p_if)_{i\in I}\)。这两个映射仅由坐标嵌入与投影给出，不依赖于有限生成元的选择；与有限生成定义域之间同态的预复合及各 \(N_i\) 上同态的后复合相容，所以同构是自然的。若 \(I=\varnothing\)，两边都是零交换群，同样的公式仍适用。
\end{proof}

这里把“每个元素的像各有有限支集”变为“存在同一个有限集控制整个像”。在整环上的有限生成挠模中，对有限多个生成元选取零化标量并取其乘积，也得到共同的非零零化标量；乘积非零依赖整环假设，而这里的支集论证对任意含幺环都成立。


\begin{chaptersummary}
模保留了向量相加和标量作用，却不再保证非零标量可逆。交换群的整数倍、正则模中的环乘法、线性算子的多项式作用都是同一组公理的实例。模同态保持加法与标量作用，其核与像因而都是子模，商掉像就得到余核，相应的同构定理也继续成立。算子模之间的同态正是保持算子作用的线性映射，模同构则对应相似关系。判断一个具体商模时，找出适当的同态及其核，往往比直接整理陪集更方便。

\ProseParagraphBreak
生成元使模的元素都能表示出来，关系说明这些表示为何不唯一。元素的零化子由使该元素为零的标量组成，全模零化子则由在整个模上作用为零的标量组成。整环上的挠元形成子模，模掉它们得到无挠商，但挠模未必具有一个共同的非零零化标量。直和只允许有限支集的元组，直积则允许任意元组；与此对应，一族坐标映射合成映入直积的映射，一族来自各因子的映射合成映出直和的映射。有限多个生成元还可把逐个元素的条件统一起来：整环上的有限生成挠模有共同的非零零化标量；任意环上的有限生成模映入直和时，其整个像只需共同的有限多个坐标。下一章将以这些具体构造为基础，系统讨论自由模、生成与关系以及正合列。
\end{chaptersummary}

\BookExercises

\begin{exercise}\label{b2:ex:02-residue-scalars}
对正整数 \(n\)，判断交换群 \(A=\mathbb Z/12\mathbb Z\) 能否具有幺性 \(\mathbb Z/n\mathbb Z\) 模结构；存在时说明其唯一性。它能否具有某个域上的向量空间结构？再判断 \((\mathbb Z/p\mathbb Z)^d\) 对素数 \(p\) 和正整数 \(d\) 的相应问题。
\end{exercise}
\begin{solution}
若有 \(\mathbb Z/n\mathbb Z\) 作用，整数作用由幺性和可加性唯一确定，而 \(\bar n=0\) 迫使 \(nA=0\)。对 \(\bar1\in A\)，这等价于 \(12\mid n\)。反过来，若 \(12\mid n\)，整数作用中 \(n\mathbb Z\) 零化全部 \(A\)，由\cref{b2:prop:annihilator-ideals}可且只能下降为所需商环作用。

若 \(A\) 是域 \(k\) 上的非零向量空间，域的特征若为零，每个非零整数标量都可逆，与 \(12A=0\) 冲突；特征若为素数 \(p\)，则 \(pA=0\)，迫使 \(12\mid p\)，也不可能。故没有这样的域结构。

对 \(B=(\mathbb Z/p\mathbb Z)^d\)，其整数零化子是 \(p\mathbb Z\)，所以恰在 \(p\mid n\) 时有唯一的 \(\mathbb Z/n\mathbb Z\) 模结构。特别地，\(B\) 有通常的 \(\mathbb F_p\) 向量空间结构，维数为 \(d\)。
\end{solution}

\begin{exercise}\label{b2:ex:02-rational-action}
设 \(A\) 是交换群。证明：\(A\) 可以成为 \(\mathbb Q\) 模，当且仅当对每个正整数 \(n\) 和每个 \(a\in A\)，方程 \(nx=a\) 都有唯一解。存在时该 \(\mathbb Q\) 模结构唯一。
\end{exercise}
\begin{solution}
若已有 \(\mathbb Q\) 作用，唯一解为 \(x=(1/n)a\)。反过来，记乘 \(n\) 的加法群自同态为 \(\mu_n\)；题设使它成为自同构，其逆也可加。对整数 \(c\) 与正整数 \(n\)，定义
\[
\left(\dfrac cn\right)a=c\mu_n^{-1}(a).
\]
若 \(c/n=d/m\)，把两种候选结果同时乘 \(mn\)，分别得到 \(mc a\) 与 \(nd a\)，由 \(mc=nd\) 相等；由于 \(\mu_{mn}\) 单射，原结果相等。因此定义不依赖分数表示。

每个 \(\mu_n^{-1}\) 都是加法群同态，故有 \((c/n)(a+b)=(c/n)a+(c/n)b\)。给两个有理数 \(c/n,d/m\)，其中 \(m,n>0\)。因为 \(\mu_{nm}^{-1}=\mu_n^{-1}\mu_m^{-1}\)，把两项改写为同分母便得到
\[
\left(\dfrac cn\right)a+\left(\dfrac dm\right)a
=\mu_{nm}^{-1}(mc a+nd a)
=\left(\dfrac{cm+dn}{nm}\right)a
=\left(\dfrac cn+\dfrac dm\right)a.
\]
再对 \((c/n)((d/m)a)\) 施加 \(\mu_{nm}\)：由于乘整数与 \(\mu_m^{-1},\mu_n^{-1}\) 可交换，所得为 \(cd a\)，与 \(((cd)/(nm))a\) 施加 \(\mu_{nm}\) 的结果相同。\(\mu_{nm}\) 单射，故
\[
\left(\dfrac cn\right)\left(\left(\dfrac dm\right)a\right)
=\left(\dfrac{cd}{nm}\right)a.
\]
\(1\) 作用为恒等映射。这些等式连同有理数作用的可加性给出全部模公理。最后，任何幺性 \(\mathbb Q\) 模结构都必须把 \(1/n\) 实现为 \(\mu_n^{-1}\)，再由整数倍决定 \(c/n\) 的作用，所以唯一。
\end{solution}

\begin{exercise}\label{b2:ex:02-intertwiner}
设 \(k\) 是域，\(V=k^2\)，\(T(a,b)=(b,0)\)，\(S=0\)。求 \(\operatorname{End}_{k[t]}(V_T)\) 的全部矩阵及其中的自同构，再求 \(V_T\to V_S\) 与 \(V_S\to V_T\) 的全部模同态。比较后两个同态空间的 \(k\) 维数，并说明它们是否包含同构。
\end{exercise}
\begin{solution}
由\cref{b2:prop:operator-module-intertwiners}，只须检验矩阵的交织等式。在标准基下，\(T\) 的矩阵为 \(J=\begin{pmatrix}0&1\\0&0\end{pmatrix}\)。设 \(U=\begin{pmatrix}a&b\\c&d\end{pmatrix}\)，计算 \(UJ=JU\) 得 \(c=0\)、\(d=a\)。所以全部自同态恰为
\[
\begin{pmatrix}a&b\\0&a\end{pmatrix},\qquad a,b\in k.
\]
其行列式为 \(a^2\)，故自同构正是 \(a\ne0\) 的那些矩阵。

从 \(V_T\) 映到 \(V_S\) 要求 \(UJ=0\)，从 \(V_S\) 映到 \(V_T\) 则要求 \(JU=0\)。因此这两个方向的同态矩阵分别为
\[
\begin{pmatrix}0&b\\0&d\end{pmatrix},\qquad
\begin{pmatrix}a&b\\0&0\end{pmatrix},\qquad a,b,d\in k.
\]
两个同态空间的 \(k\) 维数都为 \(2\)，但每个矩阵的秩至多为 \(1\)，没有可逆者。因此同态空间同维并不使这两个算子模同构。
\end{solution}

\begin{exercise}\label{b2:ex:02-hom-cyclic}
设 \(m,n\) 为正整数，\(d=\gcd(m,n)\)。求交换群
\(\operatorname{Hom}_{\mathbb Z}(\mathbb Z/m\mathbb Z,\mathbb Z/n\mathbb Z)\)，给出它的一个生成元及每个同态的公式。
\end{exercise}
\begin{solution}
同态由 \(\bar1\) 的像 \(\bar y\) 唯一决定，且必须有 \(m\bar y=0\)。反过来，这个条件保证公式 \(f(\bar a)=a\bar y\) 不依赖 \(a\) 的代表元，故给出同态。令 \(m=dm'\)、\(n=dn'\)，则 \(n\mid my\) 等价于 \(n'\mid m'y\)。由于 \(\gcd(m',n')=1\)，Bézout 等式说明这又等价于 \(n'\mid y\)。因此允许的像恰为
\[
\bar y=j\overline{n/d},\qquad j=0,1,\ldots,d-1.
\]
它们组成由 \(\overline{n/d}\) 生成的循环子群。因为 \(n=d(n/d)\)，而 \(j(n/d)\) 被 \(n\) 整除当且仅当 \(d\mid j\)，所以 \(\overline{n/d}\) 的阶恰为 \(d\)。对应同态群同构于 \(\mathbb Z/d\mathbb Z\)，一个生成元是 \(\bar a\mapsto\overline{a n/d}\)。
\end{solution}

\begin{exercise}\label{b2:ex:02-matrix-standard-module}
令 \(R=M_n(k)\) 作用在 \(M=k^n\) 的列向量上，\(n\ge1\)。证明 \(M\) 只有子模 \(0,M\)，并求 \(\operatorname{End}_R(M)\)。这里比较所得自同态环与 \(\operatorname{End}_k(k^n)\)，分别说明 \(n=1\) 与 \(n>1\) 的情形。
\end{exercise}
\begin{solution}
若子模 \(L\) 含有非零向量 \(v\)，选 \(v_j\ne0\)。对任意 \(w\in k^n\)，取第 \(j\) 列为 \(v_j^{-1}w\)、其余列为零的矩阵 \(A\)，便有 \(Av=w\)。所以 \(L=M\)。

任意 \(R\) 模自同态 \(u\) 保持标量矩阵 \(aI_n\) 的作用，故 \(k\) 线性。又有 \(u(e_j)=u(E_{jj}e_j)=E_{jj}\cdot u(e_j)\)，所以 \(u(e_j)=a_je_j\)。从 \(E_{ij}e_j=e_i\) 得
\[
a_ie_i=u(e_i)=u(E_{ij}e_j)=E_{ij}\cdot u(e_j)=a_je_i,
\]
故全部 \(a_j\) 相等。反过来，每个标量倍恒等映射都与所有矩阵作用交换。因此 \(\operatorname{End}_R(M)\cong k\)。若 \(n>1\)，它只是 \(\operatorname{End}_k(k^n)=M_n(k)\) 的真子环；\(R\) 线性要求保持全部矩阵作用，比 \(k\) 线性更强。\(n=1\) 时两者相同。
\end{solution}

\begin{exercise}\label{b2:ex:02-opposite-regular-endomorphisms}
设 \(R\) 为含幺环。证明同一个交换群上的右 \(R\) 模结构与左 \(R^{\mathrm{op}}\) 模结构一一对应。将右正则模 \(R_R\) 也看作左 \(R^{\mathrm{op}}\) 模，分别求自同态环
\[
\operatorname{End}_R({}_RR),\qquad
\operatorname{End}_{R^{\mathrm{op}}}(R_R),
\]
并说明两者的乘法次序。
\end{exercise}
\begin{solution}
反环中 \(r*s=sr\)。若 \(M\) 是右 \(R\) 模，定义左作用 \(r\cdot m=mr\)，则
\[
(r*s)\cdot m=m(sr)=(ms)r=r\cdot(s\cdot m).
\]
加法分配律及单位作用也由右模公理给出。反过来，由左 \(R^{\mathrm{op}}\) 模作用定义 \(mr=r\cdot m\)，同一等式反向验证右模公理；两个构造逐项互逆。

左 \(R\) 线性映射 \(u:R\to R\) 由 \(a=u(1_R)\) 决定，且 \(u(x)=xa\)。记它为 \(u_a\)，则
\[
(u_a\circ u_b)(x)=(xb)a=x(ba),\qquad u_a\circ u_b=u_{ba}.
\]
因此 \(a\mapsto u_a\) 给出环同构 \(R^{\mathrm{op}}\cong\operatorname{End}_R({}_RR)\)：右乘映射复合时，参数的乘法次序反转。

右 \(R\) 线性映射 \(v:R\to R\) 则由 \(a=v(1_R)\) 决定，且 \(v(x)=ax\)。此处右 \(R\) 线性与左 \(R^{\mathrm{op}}\) 线性等价，并有
\[
(v_a\circ v_b)(x)=a(bx)=(ab)x,\qquad v_a\circ v_b=v_{ab}.
\]
所以 \(R\cong\operatorname{End}_{R^{\mathrm{op}}}(R_R)\)。左乘映射复合时，参数乘法保持原次序。
\end{solution}

\begin{exercise}\label{b2:ex:02-restriction-of-scalars}
沿 \(\mathbb R\hookrightarrow\mathbb C\) 限制标量，把 \(\mathbb C\) 分别看作复向量空间和实向量空间；比较其子模与自同态环。再沿商映射 \(\mathbb Z\to\mathbb Z/n\mathbb Z\) 限制标量，其中 \(n\ge2\) 为整数，把 \(\mathbb Z/n\mathbb Z\) 分别看作自身上的模和 \(\mathbb Z\) 模，作同样的比较。
\end{exercise}
\begin{solution}
\(\mathbb C\) 作为一维复向量空间只有 \(0\) 与 \(\mathbb C\) 两个子模；限制到 \(\mathbb R\) 后，任意非零 \(z\in\mathbb C\) 还生成实直线 \(\mathbb Rz\)，所以子模为实线性子空间。复线性自同态均为乘以某个 \(a+bi\)，故其环同构于 \(\mathbb C\)。在实基 \((1,i)\) 下，这些映射对应矩阵
\[
\begin{pmatrix}a&-b\\b&a\end{pmatrix},
\]
而所有实线性自同态组成 \(M_2(\mathbb R)\)；例如复共轭是实线性的，却不是复线性的。

设 \(S=\mathbb Z/n\mathbb Z\)。因为 \(\mathbb Z\to S\) 满射，一个加法子群对整数倍封闭，当且仅当它对 \(S\) 的作用封闭，故限制标量前后的子模相同。对加法群同态 \(f:S\to S\) 及 \(s\in S\)，任选整数 \(r\) 使其像为 \(s\)，便有 \(f(sx)=f(rx)=r f(x)=sf(x)\)。因此 \(\mathbb Z\) 线性自同态与 \(S\) 线性自同态也相同，均由 \(f(\bar1)\) 决定；两种自同态环都同构于 \(S\)。这里的商映射是满射，限制标量仍保留原来的全部作用；沿 \(\mathbb R\hookrightarrow\mathbb C\) 限制标量则只保留实数的作用。
\end{solution}

\begin{exercise}\label{b2:ex:02-congruence-kernel}
对 \(\mathbb Z\) 模同态
\[
f:\mathbb Z^2\longrightarrow\mathbb Z/6\mathbb Z,\qquad
f(a,b)=\overline{2a+3b},
\]
求核、像、余核，并给出 \(\mathbb Z^2/\ker f\) 的一个生成元及其阶。
\end{exercise}
\begin{solution}
同余 \(2a+3b\equiv0\pmod6\) 分别模三与模二考察，得到 \(3\mid a\)、\(2\mid b\)，反过来这两个条件也足够。故 \(\ker f=3\mathbb Z\times2\mathbb Z\)。由于 \(f(-1,1)=\bar1\)，像为全部 \(\mathbb Z/6\mathbb Z\)，余核为零。由\cref{b2:thm:module-first-isomorphism}，\(f\) 诱导同构 \(\mathbb Z^2/\ker f\cong\mathbb Z/6\mathbb Z\)，其中 \((-1,1)+\ker f\) 对应 \(\bar1\)，故它生成商模且阶为六。
\end{solution}

\begin{exercise}\label{b2:ex:02-dual-numbers-submodules}
把 \(M=k[t]/(t^2)\) 看成 \(k[t]\) 模，令 \(x=t+(t^2)\)。列出它的全部子模，求每个非零元素的零化子，并说明哪些元素生成 \(M\)。
\end{exercise}
\begin{solution}
多项式除以 \(t^2\) 的余式使每个元素唯一写成 \(a+bx\)。若 \(a\ne0\)，则
\[
(a+bx)(a^{-1}-ba^{-2}x)=1,
\]
所以它在商环中为单位，并作为 \(k[t]\) 模生成整个 \(M\)。包含这种元素的子模只能是 \(M\)。否则子模包含在 \(kx\) 中；若含有某个 \(bx\ne0\)，常数标量 \(b^{-1}\) 使它包含 \(x\)，从而等于 \(kx\)。因此全部子模为 \(0,kx,M\)。

若 \(a\ne0\)，由上述逆元可知 \(f(t)(a+bx)=0\) 等价于 \(f(t)\in(t^2)\)，其零化子为 \((t^2)\)。若 \(a=0\)、\(b\ne0\)，则 \(f(t)bx=0\) 等价于 \(f\) 的常数项为零，零化子为 \((t)\)。生成元恰为常数项非零的元素。
\end{solution}

\begin{exercise}\label{b2:ex:02-left-ideal-annihilator}
设 \(k\) 为域，\(R=M_2(k)\times k\)，\(e_1=(1,0)^{\mathsf T}\)，并令
\[
I=\{\,(A,c)\in R\mid Ae_1=0\,\},\qquad
K=M_2(k)\times\{0\}.
\]
证明 \(I\) 是左理想但不是双边理想，\(K\) 是双边理想。给出 \(R/I\cong k^2\) 与 \(R/K\cong k\) 的具体左 \(R\) 模同构，写明右侧的作用，并计算两个商模的零化子。证明这两个商模分别不忠实，但 \((R/I)\oplus(R/K)\) 忠实。
\end{exercise}
\begin{solution}
在 \(k^2\) 上规定 \((A,c)v=Av\)，在 \(k\) 上规定 \((A,c)d=cd\)。映射 \((A,c)\mapsto Ae_1\) 是满左模同态，核为 \(I\)，故 \(I\) 是左理想，并由\cref{b2:thm:module-first-isomorphism}得到 \((A,c)+I\mapsto Ae_1\) 的同构。但 \((E_{12},0)\in I\)，而
\[
(E_{12},0)(E_{21},0)=(E_{11},0)\notin I,
\]
所以 \(I\) 不是双边理想。第二个投影 \((A,c)\mapsto c\) 是满环同态，核为双边理想 \(K\)，并给出左模同构 \((A,c)+K\mapsto c\)。

在第一个商模上，\((A,c)\) 零化全部向量当且仅当 \(A=0\)；在第二个商模上，则当且仅当 \(c=0\)。因此
\[
\operatorname{Ann}_R(R/I)=\{0\}\times k,\qquad
\operatorname{Ann}_R(R/K)=M_2(k)\times\{0\}.
\]
它们分别是 \(I,K\) 中的最大双边理想，与\cref{b2:prop:annihilator-ideals}一致。两个零化子都非零，所以两个商模都不忠实；零化它们的直和须同时属于这两个零化子，而两者交为零，故直和忠实。
\end{solution}

\begin{exercise}\label{b2:ex:02-nonfinite-submodule}
令 \(R=k[X_1,X_2,\ldots]\) 为可数无限个变量的多项式环，其中每个多项式只用有限多个变量。证明正则模 \(R\) 有限生成，但其子模 \(I=(X_1,X_2,\ldots)\) 不有限生成。
\end{exercise}
\begin{solution}
正则模由 \(1\) 生成。假设 \(I\) 可由 \(f_1,\ldots,f_t\) 生成，每个 \(f_j\) 都有零常数项，而其中出现的变量总共只有有限多个，记其指标集为 \(F\)。选 \(q\notin F\)。把 \(X_q\) 保留，其余变量均代为零，得到环同态 \(\varepsilon:R\to k[X_q]\)。每个 \(f_j\) 的像为零，所以由它们生成的理想全部映为零；但 \(X_q\in I\) 且 \(\varepsilon(X_q)=X_q\ne0\)，矛盾。这说明有限生成模的子模未必有限生成，不能直接搬用有限维向量空间的结论。
\end{solution}

\begin{exercise}\label{b2:ex:02-torsion-quotient-lattice}
设 \(R\) 为整环，\(L\le M\)。证明 \(M/L\) 无挠当且仅当对所有 \(0\ne r\in R\)、\(m\in M\)，由 \(rm\in L\) 可以推出 \(m\in L\)。对 \(M=\mathbb Z^2\)、\(L=\mathbb Z(2,4)\)，求 \(M/L\) 的挠子模及其无挠商。
\end{exercise}
\begin{solution}
等式 \(r(m+L)=0\) 等价于 \(rm\in L\)，而 \(m+L=0\) 等价于 \(m\in L\)，所以所述条件就是无挠的定义。

对具体的 \(L\)，若 \(n(a,b)=t(2,4)\)，其中 \(n\ne0\)，则 \(b=2a\)。反过来，若 \((a,b)=a(1,2)\)，则 \(2(a,b)\in L\)。故挠子模为
\[
\mathbb Z(1,2)/\mathbb Z(2,4)\cong\mathbb Z/2\mathbb Z.
\]
模掉这项后，由\cref{b2:thm:module-correspondence}得到 \(\mathbb Z^2/\mathbb Z(1,2)\)。同态 \((a,b)\mapsto b-2a\) 满射到 \(\mathbb Z\)，核为 \(\mathbb Z(1,2)\)，所以该无挠商同构于 \(\mathbb Z\)。
\end{solution}

\begin{exercise}\label{b2:ex:02-torsion-product}
取一个素数 \(p\)，令 \(M_n=\mathbb Z/p^n\mathbb Z\)、\(n\ge1\)。证明 \(\bigoplus_{n\ge1}M_n\) 是挠 \(\mathbb Z\) 模，但 \(\prod_{n\ge1}M_n\) 不是挠模。再证明任意一族无挠整环模的直积仍无挠。
\end{exercise}
\begin{solution}
直和元素只有有限个非零坐标。若这些坐标的下标均不超过 \(N\)，则 \(p^N\) 零化该元素，所以直和是挠模。直积中的 \(x=(\bar1)_{n\ge1}\) 若被非零整数 \(a\) 零化，就要求 \(p^n\mid a\) 对全部 \(n\) 成立，这与 \(a\ne0\) 矛盾。

若各 \(N_i\) 都在整环 \(R\) 上无挠，且 \(0\ne r\in R\)、\(r(n_i)=0\)，则每个 \(rn_i=0\) 都迫使 \(n_i=0\)，故整个元组为零。因此直积无挠，其子模直和当然也无挠。
\end{solution}

\begin{exercise}\label{b2:ex:02-annihilator-sum-product}
设 \((M_i)_{i\in I}\) 为一族左 \(R\) 模，证明
\[
\operatorname{Ann}_R\!\left(\bigoplus_iM_i\right)
=\operatorname{Ann}_R\!\left(\prod_iM_i\right)
=\bigcap_i\operatorname{Ann}_R(M_i).
\]
据此判断 \(\bigoplus_{n\ge1}\mathbb Z/p^n\mathbb Z\) 是否忠实，并说明它与上一题的挠性结论是否冲突。
\end{exercise}
\begin{solution}
若 \(r\) 零化每个 \(M_i\)，则逐坐标作用表明它零化直积，也零化直和。反过来，直和与直积都含有每个 \(M_i\) 的单坐标嵌入；\(r\) 零化整个模，就必定零化每个坐标模。空族时两边同为 \(R\)，等式仍成立。

具体例子的零化子为 \(\bigcap_{n\ge1}p^n\mathbb Z=0\)，所以直和忠实。它同时是挠模，这没有冲突：忠实性排除的是一个非零标量零化所有元素；挠性要求的是每个元素分别有某个非零零化标量。两处量词的次序不同。
\end{solution}

\begin{exercise}\label{b2:ex:02-finite-hom-direct-sum}
设 \(R\) 是非零交换环，\(F=R^{(\mathbb N)}\)，\(e_i\) 为第 \(i\) 个坐标向量。通过
\[
f(e_i)=\sum_{j\ge1}a_{ji}e_j,
\]
证明 \(\operatorname{End}_R(F)\) 与每列只有有限多个非零项的无限矩阵 \(A=(a_{ji})\) 一一对应；求复合的矩阵公式，并说明其中的求和为何有限。再证明 \(f(F)\) 包含在某个有限坐标子直和中，当且仅当 \(A\) 只有有限多个非零行。比较定义域为 \(R^n\) 与 \(F\) 时的区别，并用恒等映射说明列有限条件并不保证只有有限多个非零行。
\end{exercise}
\begin{solution}
每个 \(f(e_i)\) 属于直和，故第 \(i\) 列只有有限多个非零项。反过来，给定列有限矩阵，对有限展开 \(x=\sum_i x_ie_i\) 定义
\[
f(x)=\sum_i x_i\left(\sum_j a_{ji}e_j\right).
\]
外层只涉及 \(x\) 的有限支集，每个内层也是有限和，所以结果仍在直和中。唯一的坐标展开保证公式无歧义，且它保持加法与标量作用。

若 \(A,B\) 分别表示 \(f,g\)，则 \(f\circ g\) 的矩阵为 \(AB\)，其中
\[
(AB)_{ji}=\sum_{h\ge1}a_{jh}b_{hi}.
\]
对固定的 \(i\)，\(B\) 的第 \(i\) 列只有有限多个非零项，故此和有限；相应的有限多个 \(A\) 列的支集之并仍有限，所以 \(AB\) 也列有限。

若像包含在由有限指标集 \(J\) 给出的子直和中，则每个 \(f(e_i)\) 在 \(J\) 外都为零，即 \(J\) 外的行全为零。反过来，只有有限多个非零行时，每个像的支集都包含在这些行的指标集中。定义域为 \(R^n\) 时，只有有限列，各列支集之并给出共同有限的行指标集，这也正是\cref{b2:prop:finite-hom-direct-sum}中的有限生成机制。对于 \(F\) 的恒等映射，矩阵为 \((\delta_{ji})\)，每列只有一个非零项，但每行都非零；它没有共同有限的坐标范围。
\end{solution}

\begin{exercise}\label{b2:ex:02-sum-quotients}
设 \(L_i\le M_i\)。构造并证明同构
\[
\left(\bigoplus_iM_i\right)\Big/\left(\bigoplus_iL_i\right)
\cong\bigoplus_i(M_i/L_i).
\]
明确说明满射性如何使用有限支集。
\end{exercise}
\begin{solution}
逐坐标取商定义同态 \(q:\bigoplus_iM_i\to\bigoplus_i(M_i/L_i)\)。商坐标非零只可能发生在原坐标非零处，所以像确实在直和中。给定目标元组，只需对其有限个非零坐标选择代表元，其余坐标取零，便得到有限支集的原像。因此 \(q\) 满射。

元组在核中当且仅当每个坐标都属于 \(L_i\)，加上原有的有限支集条件，恰好得到 \(\ker q=\bigoplus_iL_i\)。由\cref{b2:thm:module-first-isomorphism}得到题中同构，它把 \((m_i)+\bigoplus_iL_i\) 送到 \((m_i+L_i)\)。
\end{solution}

\begin{exercise}\label{b2:ex:02-idempotent-decomposition}
设 \(e\in\operatorname{End}_R(M)\) 满足 \(e^2=e\)。证明 \(M=\operatorname{im}e\oplus\ker e\)。反过来，每个内直和分解 \(M=U\oplus V\) 是否给出这样的自同态？若 \(a\in R\) 满足 \(a^2=a\)，把这个结果用于左正则模，写出所得分解。
\end{exercise}
\begin{solution}
对任意 \(m\)，分解 \(m=e(m)+(m-e(m))\) 的首项在像中，而 \(e(m-e(m))=e(m)-e^2(m)=0\)，故第二项在核中。若 \(x=e(y)\) 又在核中，则 \(x=e(y)=e^2(y)=e(x)=0\)，所以交为零。由\cref{b2:prop:internal-direct-sum}得到所述分解。

反过来，设每个 \(m\) 唯一写成 \(u+v\)，定义 \(e(u+v)=u\)。分解的唯一性与 \(U,V\) 的子模条件保证 \(e\) 可加且 \(R\) 线性，并有 \(e^2=e\)、\(\operatorname{im}e=U\)、\(\ker e=V\)。

由\cref{b2:prop:regular-hom-endomorphism}，左正则模的自同态是右乘映射，故采用 \(e(r)=ra\)。直接检验也有 \(e(sr)=(sr)a=s(ra)=s\,e(r)\)；左乘 \(r\mapsto ar\) 则一般不能将 \(a\) 与 \(s\) 交换，未必左 \(R\) 线性。由 \(a^2=a\) 得 \(e^2=e\)。像为 \(Ra\)；若 \(ra=0\)，则 \(r=r(1-a)\)，反过来 \(r(1-a)a=0\)，所以核为 \(R(1-a)\)。因此 \({}_RR=Ra\oplus R(1-a)\)，不要求 \(a\) 位于中心。
\end{solution}

\begin{exercise}\label{b2:ex:02-product-ring-module}
设 \(R,S\) 为含幺环，\(M\) 为左 \(R\times S\) 模，令 \(e=(1_R,0)\)、\(f=(0,1_S)\)。证明 \(M=eM\oplus fM\)，其中 \(eM\) 具有自然的左 \(R\) 模结构，\(fM\) 具有自然的左 \(S\) 模结构。反过来，给一个左 \(R\) 模 \(U\) 和左 \(S\) 模 \(V\)，构造相应的 \(R\times S\) 模。
\end{exercise}
\begin{solution}
\(e,f\) 是中心元素，所以它们的作用是 \(R\times S\) 模自同态，像因而为子模。由 \(e+f=1\)、\(ef=fe=0\)，每个 \(m\) 分解为 \(em+fm\)；若 \(x\in eM\cap fM\)，则一方面 \(ex=x\)，另一方面 \(ex=0\)，故 \(x=0\)。于是得到内直和。

在 \(eM\) 上定义 \(r x=(r,0)x\)，\(1_R\) 的作用为 \(e\)，在 \(eM\) 上恰是恒等映射，故这是幺性左 \(R\) 作用。\(fM\) 上的 \(S\) 作用同样由 \(s y=(0,s)y\) 定义。反过来，在 \(U\oplus V\) 上规定
\[
(r,s)(u,v)=(ru,sv),
\]
逐坐标的模公理使它成为幺性左 \(R\times S\) 模。此时 \(e,f\) 正是两个坐标投影，其像恢复 \(U\) 与 \(V\)。
\end{solution}
